\documentclass{article}
\usepackage[margin=0.8in]{geometry}
\usepackage{amsmath,amssymb,amsthm,mathtools}
\usepackage{mathrsfs,bm,bbm}
\usepackage{graphicx,booktabs,array,multirow,tabularx,longtable}
\usepackage{enumitem}
\usepackage{xcolor}
\usepackage{algorithm,algpseudocode}
\usepackage{subcaption,tikz}
\usepackage{siunitx,cancel,accents}
\usepackage{microtype}
\usepackage[numbers,sort&compress]{natbib}
\usepackage[colorlinks=true,linkcolor=blue,citecolor=blue,urlcolor=blue]{hyperref}
\usepackage{cleveref}
\setlength{\emergencystretch}{2em}
\newtheorem{assumption}{Assumption}
\newtheorem{definition}{Definition}
\newtheorem{theorem}{Theorem}
\newtheorem{lemma}{Lemma}
\newtheorem{proposition}{Proposition}
\newtheorem{openendedquestion}{Open-ended Question}
\newtheorem{conjecture}{Conjecture}
\newcommand{\coreref}[1]{\ref{#1}}

\begin{document}

\sloppy

\section*{exp\_\allowbreak{}dte\_\allowbreak{}minimax\_\allowbreak{}comp\_\allowbreak{}frontier}

\noindent\textit{Specialization:} v1\par

\paragraph{TL;DR.} For fixed finite populations under arbitrary neighborhood interference and separate \(q=2\) target-layer moment restrictions, every supplied nonempty pair \(A,B\subseteq V_n\) with \(B\subseteq\bigcap_{i\in A}\widetilde N_i(G)\) yields an explicit four-atom Bayes certificate \(R_2(G)\geq\min\{|A|,|B|\}/n\). Maximization gives \(R_2(G)\geq\beta_{\mathrm{DTE}}(G)/n\), which contains the maximum-clique and clique-partition certificates and is rate-matched by the published spectral upper functional on complete-bipartite-plus-isolates networks. On equal block-clique graphs, an explicit mode SDP gives the exact invariant-linear value and a realizable unrestricted upper endpoint, while a finite posterior prior gives a computable lower endpoint; unrestricted equality is asserted only when those endpoints verify equal, and it holds unconditionally for a single clique. A randomized polynomial-time constant-factor two-sided saddle on arbitrary DTE conflict graphs remains open.

\section{Project justification}

\paragraph{Gap.} Published \(q=2\) DTE theory supplies conflict-graph and spectral-risk comparisons, and published Conflict Graph Design supplies realizable upper-risk procedures, but it does not give a matching finite-prior certificate against all assignment laws and measurable estimators or a two-sided polynomial-time approximation to the unrestricted minimax value.

\paragraph{Niche.} The unresolved object is the conflict-constrained two-layer decision game: observing an all-control target and observing an isolated-treatment target can be mutually exclusive across large cross-layer supports, while the investigator must jointly choose a realizable randomization and an estimator against separate layerwise Euclidean balls.

\paragraph{Fill.} The paper supplies the explicit balanced-biclique four-atom lower certificate, its clique and clique-partition corollaries, and a complete-bipartite family on which it matches the published spectral upper order. For equal block cliques it also supplies an exact invariant-linear mode SDP, a realizable upper procedure, and a finite posterior lower certificate. These results expose, but do not close, the missing general polynomial-time two-sided saddle problem.

\section{Related work}

Kandiros, Harshaw, and Sävje (2026) provide the published \(q=2\) arbitrary-neighborhood-interference class, graph-risk monotonicity, the labelled conflict-graph characterization, and general graph-functional risk bounds. Kandiros, Pipis, Daskalakis, and Harshaw (2024) provide an implementable Conflict Graph Design, a modified Horvitz--Thompson estimator, and computable conservative upper guarantees, but not the matching all-procedure finite-prior lower certificate sought here. Fatemi and Zheleva (2023) provide CauseIS, whose maximal-independent-set design motivates instance-specific DTE benchmarking; the balanced-biclique certificate changes that benchmark on complete-bipartite networks where the clique witness remains constant. Gabler (1990) and Rinott (2009) inform the invariant and joint design-estimation decision framework. The present result is a stronger lower certificate and a symmetric-family computable sandwich, not a proof of the general computational minimax frontier.

\section{Setup}

Target (estimand / structural parameter): \(\tau_{\mathrm{DTE}}(Y)=n^{-1}\sum_{i\in V_n}\{Y_i(e_{i,1})-Y_i(e_{i,0})\}\).
Primary functional / estimator / diagnostic: For every known graph \(G\) and every supplied feasible nonempty ordered pair \(A,B\subseteq V_n\) satisfying \(B\subseteq\bigcap_{i\in A}\widetilde N_i(G)\), the directly evaluated four-atom lower certificate is
\[
R_2(G)\geq\frac{\min\{|A|,|B|\}}{n};
\]
maximizing this finite witness functional yields
\[
R_2(G)\geq\frac{\beta_{\mathrm{DTE}}(G)}n\geq\frac{\omega(G)}n\geq\sum_{B\in\mathcal P}\left(\frac{|B|}{n}\right)^2.
\]
Verification of a supplied pair is direct, but no polynomial-time algorithm for maximizing \(\beta_{\mathrm{DTE}}(G)\) or \(\omega(G)\) is claimed.  On an equal block-clique graph the clique corollary gives \(1/b\), and the computable two-sided functional remains
\[
\max\left\{\frac1b,\min_{0\leq c\leq b}B_c(\Pi_{b,s}^{\star})\right\}
\leq R_2(G_{b,s})\leq V_{b,s}
=\inf_{p,A,D}\sup_{u\in\Theta_{b,s}}\mathbb E\!\left[\left\{\delta_{p,A,D}-\tau_{\mathrm{DTE}}(Y)\right\}^2\right].
\]
Equality \(R_2(G_{b,s})=V_{b,s}\) is certified only when the displayed endpoints agree; for \(b=1\), both equal one.  For \(G_n=K_{m_n,m_n}\sqcup(n-2m_n)K_1\), where \(m_n=\lfloor\rho n^{\gamma}\rfloor\), \(\rho\in(0,1/2)\), and \(\gamma\in(0,1]\), the shore witness and cited spectral upper functional give \(R_2(G_n)\asymp(m_n+1)/n=\Theta(n^{\gamma-1})\), while the clique certificate is only \(2/n\) once \(m_n\geq1\).  This is a lower-certificate calibration, not a general realizable upper design or a polynomial-support two-sided saddle..
\begin{itemize}
  \item \texttt{b} --- \texttt{number\_\allowbreak{}of\_\allowbreak{}blocks}; \(\{1,2,\ldots\}\); \texttt{family\_\allowbreak{}index}
  \par\smallskip\noindent\emph{Definition.} \texttt{Number of equal clique blocks.}
  \item \texttt{s} --- \texttt{block\_\allowbreak{}size}; \(\{2,3,\ldots\}\); \texttt{family\_\allowbreak{}index}
  \par\smallskip\noindent\emph{Definition.} \texttt{Number of units in each clique block.}
  \item \texttt{n} --- \texttt{population\_\allowbreak{}size}; \(\mathbb N\); \texttt{instance\_\allowbreak{}size}
  \par\smallskip\noindent\emph{Definition.} \(n=bs\).
  \item \texttt{V\_\allowbreak{}n} --- \texttt{finite\_\allowbreak{}population}; \(2^{\mathbb N}\); \texttt{population}
  \par\smallskip\noindent\emph{Definition.} Fixed population \(V_n=\{1,\ldots,n\}\).
  \item \texttt{r} --- \texttt{block\_\allowbreak{}index}; \(\{1,\ldots,b\}\); \texttt{index}
  \par\smallskip\noindent\emph{Definition.} \texttt{Generic clique-\allowbreak{}block index.}
  \item \texttt{j} --- \texttt{within\_\allowbreak{}block\_\allowbreak{}index}; \(\{1,\ldots,s\}\); \texttt{index}
  \par\smallskip\noindent\emph{Definition.} \texttt{Generic label within a clique block.}
  \item \texttt{B\_\allowbreak{}r} --- \texttt{clique\_\allowbreak{}block}; \(2^{V_n}\); \texttt{network\_\allowbreak{}structure}
  \par\smallskip\noindent\emph{Definition.} Block \(r\) in a fixed labeled partition of \(V_n\), with \(|B_r|=s\).
  \item \texttt{\textbackslash{}\allowbreak{}mathcal B\_\allowbreak{}\{b,\allowbreak{}s\}\allowbreak{}} --- \texttt{equal\_\allowbreak{}block\_\allowbreak{}partition}; \(2^{2^{V_n}}\); \texttt{network\_\allowbreak{}structure}
  \par\smallskip\noindent\emph{Definition.} \(\mathcal B_{b,s}=\{B_1,\ldots,B_b\}\) is a partition of \(V_n\) into equal blocks.
  \item \texttt{G\_\allowbreak{}\{b,\allowbreak{}s\}\allowbreak{}} --- \texttt{block\_\allowbreak{}clique\_\allowbreak{}network}; \(2^{\binom{V_n}{2}}\); \texttt{design\_\allowbreak{}input}
  \par\smallskip\noindent\emph{Definition.} Disjoint union of the complete graphs induced by the blocks in \(\mathcal B_{b,s}\).
  \item \texttt{G} --- \texttt{known\_\allowbreak{}interference\_\allowbreak{}network}; \(2^{\binom{V_n}{2}}\); \texttt{open\_\allowbreak{}question\_\allowbreak{}input}
  \par\smallskip\noindent\emph{Definition.} Generic known simple interference graph on \(V_n\).
  \item \texttt{H\_\allowbreak{}\{\textbackslash{}\allowbreak{}mathrm\{DTE\}\allowbreak{}\}\allowbreak{}(G)} --- \texttt{general\_\allowbreak{}two\_\allowbreak{}layer\_\allowbreak{}dte\_\allowbreak{}conflict\_\allowbreak{}graph}; \(2^{\binom{V_n\times\{0,1\}}{2}}\); \texttt{open\_\allowbreak{}question\_\allowbreak{}input}
  \par\smallskip\noindent\emph{Definition.} Two-layer exposure conflict graph for the DTE on \(G\).
  \item \texttt{\textbackslash{}\allowbreak{}Theta\_\allowbreak{}2(G)} --- \texttt{general\_\allowbreak{}separate\_\allowbreak{}layer\_\allowbreak{}second\_\allowbreak{}moment\_\allowbreak{}set}; \(2^{\mathbb R^{V_n\times\{0,1\}}}\); \texttt{open\_\allowbreak{}question\_\allowbreak{}parameter\_\allowbreak{}space}
  \par\smallskip\noindent\emph{Definition.} Two-layer target-outcome vectors on \(G\) satisfying a normalized Euclidean-ball constraint in each layer.
  \item \texttt{i} --- \texttt{unit\_\allowbreak{}index}; \(V_n\); \texttt{index}
  \par\smallskip\noindent\emph{Definition.} \texttt{Generic experimental-\allowbreak{}unit index.}
  \item \texttt{\textbackslash{}\allowbreak{}widetilde N\_\allowbreak{}i} --- \texttt{closed\_\allowbreak{}neighborhood}; \(2^{V_n}\); \texttt{interference\_\allowbreak{}support}
  \par\smallskip\noindent\emph{Definition.} Closed neighborhood of unit \(i\) in \(G_{b,s}\); it is the unique block containing \(i\).
  \item \texttt{\textbackslash{}\allowbreak{}mathcal Z\_\allowbreak{}n} --- \texttt{assignment\_\allowbreak{}space}; \(2^{V_n}\); \texttt{assignment\_\allowbreak{}space}
  \par\smallskip\noindent\emph{Definition.} \(\mathcal Z_n=\{0,1\}^{V_n}\).
  \item \texttt{z} --- \texttt{assignment}; \(\mathcal Z_n\); \texttt{assignment\_\allowbreak{}argument}
  \par\smallskip\noindent\emph{Definition.} \texttt{Generic binary treatment assignment.}
  \item \texttt{z\textasciicircum{}\{\textbackslash{}\allowbreak{}prime\}\allowbreak{}} --- \texttt{assignment}; \(\mathcal Z_n\); \texttt{assignment\_\allowbreak{}argument}
  \par\smallskip\noindent\emph{Definition.} \texttt{Second generic treatment assignment.}
  \item \texttt{Y} --- \texttt{finite\_\allowbreak{}population\_\allowbreak{}potential\_\allowbreak{}outcome\_\allowbreak{}schedule}; \(\prod_{i\in V_n}\mathbb R^{\mathcal Z_n}\); \texttt{fixed\_\allowbreak{}finite\_\allowbreak{}population}
  \par\smallskip\noindent\emph{Definition.} Fixed schedule \(Y=(Y_i(z):i\in V_n,z\in\mathcal Z_n)\).
  \item \texttt{k} --- \texttt{exposure\_\allowbreak{}layer\_\allowbreak{}index}; \(\{0,1\}\); \texttt{index}
  \par\smallskip\noindent\emph{Definition.} \texttt{DTE exposure-\allowbreak{}layer index.}
  \item \texttt{e\_\allowbreak{}\{i,\allowbreak{}k\}\allowbreak{}} --- \texttt{dte\_\allowbreak{}target\_\allowbreak{}exposure}; \(\{0,1\}^{\widetilde N_i}\); \texttt{target\_\allowbreak{}exposure}
  \par\smallskip\noindent\emph{Definition.} \(e_{i,0}\) is all control on \(\widetilde N_i\), while \(e_{i,1}\) treats only \(i\).
  \item \texttt{u\_\allowbreak{}\{i,\allowbreak{}k\}\allowbreak{}} --- \texttt{target\_\allowbreak{}exposure\_\allowbreak{}outcome}; \(\mathbb R\); \texttt{causal\_\allowbreak{}primitive}
  \par\smallskip\noindent\emph{Definition.} \(u_{i,k}=Y_i(e_{i,k})\).
  \item \texttt{u} --- \texttt{two\_\allowbreak{}layer\_\allowbreak{}target\_\allowbreak{}outcome\_\allowbreak{}vector}; \(\mathbb R^{V_n\times\{0,1\}}\); \texttt{game\_\allowbreak{}parameter}
  \par\smallskip\noindent\emph{Definition.} Vector of all \(u_{i,k}\).
  \item \texttt{\textbackslash{}\allowbreak{}tau\_\allowbreak{}\{\textbackslash{}\allowbreak{}mathrm\{DTE\}\allowbreak{}\}\allowbreak{}(Y)} --- \texttt{finite\_\allowbreak{}population\_\allowbreak{}direct\_\allowbreak{}treatment\_\allowbreak{}effect}; \(\mathbb R\); \texttt{causal\_\allowbreak{}estimand}
  \par\smallskip\noindent\emph{Definition.} \(\tau_{\mathrm{DTE}}(Y)=n^{-1}\sum_{i\in V_n}(u_{i,1}-u_{i,0})\).
  \item \texttt{\textbackslash{}\allowbreak{}Theta\_\allowbreak{}\{b,\allowbreak{}s\}\allowbreak{}} --- \texttt{separate\_\allowbreak{}layer\_\allowbreak{}second\_\allowbreak{}moment\_\allowbreak{}set}; \(2^{\mathbb R^{V_n\times\{0,1\}}}\); \texttt{parameter\_\allowbreak{}space}
  \par\smallskip\noindent\emph{Definition.} \texttt{Product of the two normalized Euclidean balls for the control and treated target-\allowbreak{}outcome layers.}
  \item \texttt{\textbackslash{}\allowbreak{}mathcal M\_\allowbreak{}2(G\_\allowbreak{}\{b,\allowbreak{}s\}\allowbreak{})} --- \texttt{published\_\allowbreak{}second\_\allowbreak{}moment\_\allowbreak{}ani\_\allowbreak{}class}; \(2^{\prod_{i\in V_n}\mathbb R^{\mathcal Z_n}}\); \texttt{adversarial\_\allowbreak{}schedule\_\allowbreak{}class}
  \par\smallskip\noindent\emph{Definition.} Published \(q=2\) arbitrary-neighborhood-interference class specialized to \(G_{b,s}\).
  \item \texttt{\textbackslash{}\allowbreak{}mathcal D\_\allowbreak{}n} --- \texttt{assignment\_\allowbreak{}law\_\allowbreak{}space}; \(\Delta(\mathcal Z_n)\); \texttt{design\_\allowbreak{}space}
  \par\smallskip\noindent\emph{Definition.} \texttt{All randomization laws on binary assignments.}
  \item \texttt{\textbackslash{}\allowbreak{}Gamma\_\allowbreak{}n} --- \texttt{estimator\_\allowbreak{}space}; \(\{\widehat\tau:\mathcal Z_n\times\mathbb R^n\to\mathbb R\}\); \texttt{estimator\_\allowbreak{}space}
  \par\smallskip\noindent\emph{Definition.} \texttt{All measurable assignment-\allowbreak{}and-\allowbreak{}outcome estimators.}
  \item \texttt{R\_\allowbreak{}2(G\_\allowbreak{}\{b,\allowbreak{}s\}\allowbreak{})} --- \texttt{unrestricted\_\allowbreak{}dte\_\allowbreak{}minimax\_\allowbreak{}risk}; \([0,\infty)\); \texttt{minimax\_\allowbreak{}value}
  \par\smallskip\noindent\emph{Definition.} Worst-case mean squared error minimized jointly over \(\mathcal D_n\) and \(\Gamma_n\).
  \item \texttt{H\_\allowbreak{}\{b,\allowbreak{}s\}\allowbreak{}} --- \texttt{two\_\allowbreak{}layer\_\allowbreak{}dte\_\allowbreak{}conflict\_\allowbreak{}graph}; \(2^{\binom{V_n\times\{0,1\}}{2}}\); \texttt{calibration\_\allowbreak{}graph}
  \par\smallskip\noindent\emph{Definition.} DTE exposure-level conflict graph induced by \(G_{b,s}\).
  \item \texttt{c} --- \texttt{all\_\allowbreak{}control\_\allowbreak{}block\_\allowbreak{}count}; \(\{0,\ldots,b\}\); \texttt{orbit\_\allowbreak{}index}
  \par\smallskip\noindent\emph{Definition.} \texttt{Number of blocks assigned the all-\allowbreak{}control exposure pattern.}
  \item \texttt{\textbackslash{}\allowbreak{}mathcal S\_\allowbreak{}c} --- \texttt{maximal\_\allowbreak{}exposure\_\allowbreak{}pattern\_\allowbreak{}orbit}; \(2^{2^{V_n\times\{0,1\}}}\); \texttt{observation\_\allowbreak{}orbit}
  \par\smallskip\noindent\emph{Definition.} Orbit of maximal compatible exposure sets with exactly \(c\) all-control blocks and one isolated-treated label in each remaining block.
  \item \texttt{S\_\allowbreak{}c} --- \texttt{canonical\_\allowbreak{}maximal\_\allowbreak{}exposure\_\allowbreak{}pattern}; \(\mathcal S_c\); \texttt{posterior\_\allowbreak{}observation\_\allowbreak{}pattern}
  \par\smallskip\noindent\emph{Definition.} Fixed canonical representative of orbit \(\mathcal S_c\).
  \item \texttt{C} --- \texttt{all\_\allowbreak{}control\_\allowbreak{}block\_\allowbreak{}subset}; \(2^{\{1,\ldots,b\}}\); \texttt{assignment\_\allowbreak{}component}
  \par\smallskip\noindent\emph{Definition.} Set of all-control blocks, with \(|C|=c\).
  \item \texttt{J\_\allowbreak{}r} --- \texttt{singleton\_\allowbreak{}treated\_\allowbreak{}label}; \(\{1,\ldots,s\}\); \texttt{assignment\_\allowbreak{}component}
  \par\smallskip\noindent\emph{Definition.} Treated label selected in block \(r\notin C\).
  \item \texttt{p} --- \texttt{orbit\_\allowbreak{}law}; \(\Delta(\{0,\ldots,b\})\); \texttt{invariant\_\allowbreak{}design\_\allowbreak{}parameter}
  \par\smallskip\noindent\emph{Definition.} Law \(p=(p_0,\ldots,p_b)\) of the all-control block count.
  \item \texttt{A\_\allowbreak{}c} --- \texttt{treated\_\allowbreak{}observation\_\allowbreak{}coefficient}; \(\mathbb R\); \texttt{invariant\_\allowbreak{}estimator\_\allowbreak{}parameter}
  \par\smallskip\noindent\emph{Definition.} Common coefficient on each observed isolated-treated outcome when the orbit index is \(c\).
  \item \texttt{D\_\allowbreak{}c} --- \texttt{control\_\allowbreak{}observation\_\allowbreak{}coefficient}; \(\mathbb R\); \texttt{invariant\_\allowbreak{}estimator\_\allowbreak{}parameter}
  \par\smallskip\noindent\emph{Definition.} Common coefficient on each observed control outcome when the orbit index is \(c\).
  \item \texttt{\textbackslash{}\allowbreak{}delta\_\allowbreak{}\{p,\allowbreak{}A,\allowbreak{}D\}\allowbreak{}} --- \texttt{invariant\_\allowbreak{}design\_\allowbreak{}estimator\_\allowbreak{}procedure}; \(\mathcal D_n\times\Gamma_n\); \texttt{candidate\_\allowbreak{}saddle\_\allowbreak{}action}
  \par\smallskip\noindent\emph{Definition.} Block- and unit-invariant randomization and estimator indexed by \(p\), \(A_c\), and \(D_c\).
  \item \texttt{W\_\allowbreak{}c} --- \texttt{conditional\_\allowbreak{}coefficient\_\allowbreak{}moment\_\allowbreak{}matrix}; \(\mathbb S_+^3\); \texttt{sdp\_\allowbreak{}primal\_\allowbreak{}variable}
  \par\smallskip\noindent\emph{Definition.} Convexified moment matrix with rank-one realization \(W_c=p_c(1,A_c,D_c)^{\mathsf T}(1,A_c,D_c)\).
  \item \texttt{\textbackslash{}\allowbreak{}alpha} --- \texttt{coefficient\_\allowbreak{}moment\_\allowbreak{}index}; \(\{0,1,2\}\); \texttt{index}
  \par\smallskip\noindent\emph{Definition.} \texttt{First coordinate index in a coefficient moment matrix.}
  \item \texttt{\textbackslash{}\allowbreak{}beta} --- \texttt{coefficient\_\allowbreak{}moment\_\allowbreak{}index}; \(\{0,1,2\}\); \texttt{index}
  \par\smallskip\noindent\emph{Definition.} \texttt{Second coordinate index in a coefficient moment matrix.}
  \item \texttt{w\textasciicircum{}c\_\allowbreak{}\{\textbackslash{}\allowbreak{}alpha\textbackslash{}\allowbreak{}beta\}\allowbreak{}} --- \texttt{coefficient\_\allowbreak{}moment\_\allowbreak{}entry}; \(\mathbb R\); \texttt{sdp\_\allowbreak{}coordinate}
  \par\smallskip\noindent\emph{Definition.} Entry \((\alpha,\beta)\) of \(W_c\), ordered as \((1,A_c,D_c)\).
  \item \texttt{K\_\allowbreak{}0(W)} --- \texttt{constant\_\allowbreak{}block\_\allowbreak{}mean\_\allowbreak{}risk\_\allowbreak{}matrix}; \(\mathbb S_+^2\); \texttt{mode\_\allowbreak{}risk}
  \par\smallskip\noindent\emph{Definition.} \texttt{Two-\allowbreak{}layer risk matrix on the constant block-\allowbreak{}mean mode.}
  \item \texttt{K\_\allowbreak{}\{\textbackslash{}\allowbreak{}mathrm B\}\allowbreak{}(W)} --- \texttt{between\_\allowbreak{}block\_\allowbreak{}risk\_\allowbreak{}matrix}; \(\mathbb S_+^2\); \texttt{mode\_\allowbreak{}risk}
  \par\smallskip\noindent\emph{Definition.} \texttt{Two-\allowbreak{}layer risk matrix on every between-\allowbreak{}block zero-\allowbreak{}sum mode.}
  \item \texttt{k\_\allowbreak{}\{\textbackslash{}\allowbreak{}mathrm W\}\allowbreak{}(W)} --- \texttt{within\_\allowbreak{}block\_\allowbreak{}treated\_\allowbreak{}contrast\_\allowbreak{}risk}; \([0,\infty)\); \texttt{mode\_\allowbreak{}risk}
  \par\smallskip\noindent\emph{Definition.} \texttt{Risk eigenvalue on the treated within-\allowbreak{}block contrast modes.}
  \item \texttt{\textbackslash{}\allowbreak{}lambda\_\allowbreak{}0} --- \texttt{control\_\allowbreak{}layer\_\allowbreak{}risk\_\allowbreak{}multiplier}; \([0,\infty)\); \texttt{sdp\_\allowbreak{}epigraph\_\allowbreak{}variable}
  \par\smallskip\noindent\emph{Definition.} \texttt{Dual norm multiplier for the control-\allowbreak{}layer Euclidean ball.}
  \item \texttt{\textbackslash{}\allowbreak{}lambda\_\allowbreak{}1} --- \texttt{treated\_\allowbreak{}layer\_\allowbreak{}risk\_\allowbreak{}multiplier}; \([0,\infty)\); \texttt{sdp\_\allowbreak{}epigraph\_\allowbreak{}variable}
  \par\smallskip\noindent\emph{Definition.} \texttt{Dual norm multiplier for the treated-\allowbreak{}layer Euclidean ball.}
  \item \texttt{V\_\allowbreak{}\{b,\allowbreak{}s\}\allowbreak{}} --- \texttt{block\_\allowbreak{}clique\_\allowbreak{}invariant\_\allowbreak{}linear\_\allowbreak{}value}; \([0,\infty)\); \texttt{invariant\_\allowbreak{}linear\_\allowbreak{}value\_\allowbreak{}and\_\allowbreak{}unrestricted\_\allowbreak{}upper\_\allowbreak{}certificate}
  \par\smallskip\noindent\emph{Definition.} \texttt{Optimal value of the explicit mode semidefinite program;\allowbreak{} it is the exact minimax value within the invariant linear class and a computable unrestricted upper-\allowbreak{}risk endpoint.}
  \item \texttt{\textbackslash{}\allowbreak{}mathfrak G\_\allowbreak{}\{b,\allowbreak{}s\}\allowbreak{}} --- \texttt{block\_\allowbreak{}unit\_\allowbreak{}symmetry\_\allowbreak{}group}; \(\operatorname{Sym}(V_n)\times\{\pm1\}\); \texttt{invariance\_\allowbreak{}group}
  \par\smallskip\noindent\emph{Definition.} \texttt{Wreath-\allowbreak{}product permutations of blocks and labels within blocks,\allowbreak{} augmented by simultaneous sign reversal of both outcome layers.}
  \item \texttt{g\_\allowbreak{}0} --- \texttt{constant\_\allowbreak{}mode\_\allowbreak{}vector}; \(\mathbb R^{V_n}\); \texttt{certificate\_\allowbreak{}mode\_\allowbreak{}basis}
  \par\smallskip\noindent\emph{Definition.} \(g_0=n^{-1/2}\mathbf 1_{V_n}\).
  \item \texttt{g\_\allowbreak{}\{\textbackslash{}\allowbreak{}mathrm B\}\allowbreak{}} --- \texttt{canonical\_\allowbreak{}between\_\allowbreak{}block\_\allowbreak{}vector}; \(\mathbb R^{V_n}\); \texttt{certificate\_\allowbreak{}mode\_\allowbreak{}basis}
  \par\smallskip\noindent\emph{Definition.} For \(b\ge2\), \(g_{\mathrm B}=(2s)^{-1/2}(\mathbf 1_{B_1}-\mathbf 1_{B_2})\); for \(b=1\), set \(g_{\mathrm B}=0\).
  \item \texttt{g\_\allowbreak{}\{\textbackslash{}\allowbreak{}mathrm W\}\allowbreak{}} --- \texttt{canonical\_\allowbreak{}within\_\allowbreak{}block\_\allowbreak{}vector}; \(\mathbb R^{V_n}\); \texttt{certificate\_\allowbreak{}mode\_\allowbreak{}basis}
  \par\smallskip\noindent\emph{Definition.} \(g_{\mathrm W}=2^{-1/2}(\mathbf e_{(1,1)}-\mathbf e_{(1,2)})\), using the first two labels of block \(B_1\).
  \item \texttt{P\_\allowbreak{}c} --- \texttt{revealed\_\allowbreak{}coordinate\_\allowbreak{}projection}; \(\mathbb R^{2n\times2n}\); \texttt{posterior\_\allowbreak{}collision\_\allowbreak{}operator}
  \par\smallskip\noindent\emph{Definition.} Orthogonal coordinate projection onto the target outcomes revealed by the canonical maximal exposure pattern \(S_c\).
  \item \texttt{\textbackslash{}\allowbreak{}mathcal A\_\allowbreak{}\{b,\allowbreak{}s\}\allowbreak{}\textasciicircum{}\{\textbackslash{}\allowbreak{}star\}\allowbreak{}} --- \texttt{finite\_\allowbreak{}candidate\_\allowbreak{}orbit\_\allowbreak{}atom\_\allowbreak{}dictionary}; \(2^{\Theta_{b,s}}\); \texttt{candidate\_\allowbreak{}posterior\_\allowbreak{}collision\_\allowbreak{}support}
  \par\smallskip\noindent\emph{Definition.} Finite candidate atom dictionary produced from an optimal mode-SDP primal--dual pair by the posterior-collision lift and closure under \(\mathfrak G_{b,s}\).
  \item \texttt{L\_\allowbreak{}\{b,\allowbreak{}s\}\allowbreak{}} --- \texttt{certificate\_\allowbreak{}support\_\allowbreak{}size}; \(\mathbb N\); \texttt{certificate\_\allowbreak{}size}
  \par\smallskip\noindent\emph{Definition.} Cardinality of \(\mathcal A_{b,s}^{\star}\).
  \item \texttt{h} --- \texttt{certificate\_\allowbreak{}atom\_\allowbreak{}index}; \(\{1,\ldots,L_{b,s}\}\); \texttt{index}
  \par\smallskip\noindent\emph{Definition.} \texttt{Index of a lifted symmetry-\allowbreak{}orbit atom.}
  \item \texttt{u\textasciicircum{}\{(h)\}\allowbreak{}} --- \texttt{certificate\_\allowbreak{}atom}; \(\Theta_{b,s}\); \texttt{least\_\allowbreak{}favourable\_\allowbreak{}atom}
  \par\smallskip\noindent\emph{Definition.} Atom \(h\) in \(\mathcal A_{b,s}^{\star}\).
  \item \texttt{q\_\allowbreak{}h} --- \texttt{certificate\_\allowbreak{}atom\_\allowbreak{}weight}; \([0,1]\); \texttt{least\_\allowbreak{}favourable\_\allowbreak{}weight}
  \par\smallskip\noindent\emph{Definition.} Weight on atom \(u^{(h)}\), with weights constant on every \(\mathfrak G_{b,s}\)-orbit.
  \item \texttt{\textbackslash{}\allowbreak{}Pi\_\allowbreak{}\{b,\allowbreak{}s\}\allowbreak{}\textasciicircum{}\{\textbackslash{}\allowbreak{}star\}\allowbreak{}} --- \texttt{finite\_\allowbreak{}maximin\_\allowbreak{}posterior\_\allowbreak{}lower\_\allowbreak{}certificate}; \(\Delta(\Theta_{b,s})\); \texttt{finite\_\allowbreak{}posterior\_\allowbreak{}lower\_\allowbreak{}certificate}
  \par\smallskip\noindent\emph{Definition.} \(\Pi_{b,s}^{\star}=\sum_{h=1}^{L_{b,s}}q_h\delta_{u^{(h)}}\) is the finite prior selected to maximize the minimum orbitwise posterior risk over the candidate dictionary.
  \item \texttt{B\_\allowbreak{}c(\textbackslash{}\allowbreak{}Pi)} --- \texttt{orbitwise\_\allowbreak{}posterior\_\allowbreak{}risk}; \([0,\infty)\); \texttt{dual\_\allowbreak{}certificate\_\allowbreak{}constraint}
  \par\smallskip\noindent\emph{Definition.} Bayes residual mean squared error after observing any maximal exposure pattern in \(\mathcal S_c\) under finite prior \(\Pi\).
  \item \texttt{\textbackslash{}\allowbreak{}mathcal Q} --- \texttt{posterior\_\allowbreak{}atom\_\allowbreak{}cell}; \(2^{\{1,\ldots,L_{b,s}\}}\); \texttt{posterior\_\allowbreak{}partition}
  \par\smallskip\noindent\emph{Definition.} Cell of prior atoms with identical revealed subvectors under \(S_c\).
  \item \texttt{m\_\allowbreak{}\{\textbackslash{}\allowbreak{}mathcal Q\}\allowbreak{}} --- \texttt{posterior\_\allowbreak{}cell\_\allowbreak{}mean}; \(\mathbb R\); \texttt{bayes\_\allowbreak{}action}
  \par\smallskip\noindent\emph{Definition.} Posterior mean of the DTE over cell \(\mathcal Q\).
  \item \texttt{\textbackslash{}\allowbreak{}mathcal C\textasciicircum{}\{\textbackslash{}\allowbreak{}star\}\allowbreak{}} --- \texttt{active\_\allowbreak{}orbit\_\allowbreak{}set}; \(2^{\{0,\ldots,b\}}\); \texttt{complementary\_\allowbreak{}slackness\_\allowbreak{}set}
  \par\smallskip\noindent\emph{Definition.} \(\mathcal C^{\star}=\{c:p_c^{\star}>0\}\) for the rank-one optimal mode-SDP solution.
  \item \texttt{\textbackslash{}\allowbreak{}widetilde N\_\allowbreak{}i(G)} --- \texttt{general\_\allowbreak{}graph\_\allowbreak{}closed\_\allowbreak{}neighborhood}; \(2^{V_n}\); \texttt{general\_\allowbreak{}interference\_\allowbreak{}support}
  \par\smallskip\noindent\emph{Definition.} Closed neighborhood \(\widetilde N_i(G)=\{i\}\cup\{j\in V_n:\{i,j\}\in E(G)\}\).
  \item \texttt{\textbackslash{}\allowbreak{}tau\_\allowbreak{}\{\textbackslash{}\allowbreak{}mathrm\{DTE\}\allowbreak{}\}\allowbreak{}\textasciicircum{}G(Y)} --- \texttt{general\_\allowbreak{}graph\_\allowbreak{}finite\_\allowbreak{}population\_\allowbreak{}direct\_\allowbreak{}treatment\_\allowbreak{}effect}; \(\mathbb R\); \texttt{general\_\allowbreak{}causal\_\allowbreak{}estimand}
  \par\smallskip\noindent\emph{Definition.} Graph-indexed fixed-population DTE \(\tau_{\mathrm{DTE}}^G(Y)=n^{-1}\sum_{i\in V_n}\{Y_i(\mathbf e_i)-Y_i(\mathbf 0)\}\).
  \item \texttt{\textbackslash{}\allowbreak{}mathcal M\_\allowbreak{}2(G)} --- \texttt{general\_\allowbreak{}published\_\allowbreak{}second\_\allowbreak{}moment\_\allowbreak{}ani\_\allowbreak{}class}; \(2^{\prod_{i\in V_n}\mathbb R^{\mathcal Z_n}}\); \texttt{general\_\allowbreak{}adversarial\_\allowbreak{}schedule\_\allowbreak{}class}
  \par\smallskip\noindent\emph{Definition.} Published \(q=2\) arbitrary-neighborhood-interference schedule class on \(G\) with separate target-layer second-moment restrictions.
  \item \texttt{R\_\allowbreak{}2(G)} --- \texttt{general\_\allowbreak{}unrestricted\_\allowbreak{}dte\_\allowbreak{}minimax\_\allowbreak{}risk}; \([0,\infty)\); \texttt{general\_\allowbreak{}minimax\_\allowbreak{}value}
  \par\smallskip\noindent\emph{Definition.} Worst-case DTE mean squared error minimized jointly over all assignment laws and measurable assignment-and-observed-outcome estimators on \(G\).
  \item \texttt{\textbackslash{}\allowbreak{}mathcal P} --- \texttt{clique\_\allowbreak{}partition}; \(2^{2^{V_n}}\); \texttt{lower\_\allowbreak{}certificate\_\allowbreak{}partition}
  \par\smallskip\noindent\emph{Definition.} A partition of \(V_n\) into nonempty vertex sets that induce cliques of \(G\).
  \item \texttt{B} --- \texttt{generic\_\allowbreak{}clique\_\allowbreak{}partition\_\allowbreak{}block}; \(2^{V_n}\); \texttt{clique\_\allowbreak{}partition\_\allowbreak{}index}
  \par\smallskip\noindent\emph{Definition.} Generic nonempty block in the clique partition \(\mathcal P\).
  \item \texttt{H\_\allowbreak{}\{\textbackslash{}\allowbreak{}mathcal P\}\allowbreak{}} --- \texttt{spanning\_\allowbreak{}disjoint\_\allowbreak{}union\_\allowbreak{}of\_\allowbreak{}partition\_\allowbreak{}cliques}; \(2^{\binom{V_n}{2}}\); \texttt{lower\_\allowbreak{}certificate\_\allowbreak{}subgraph}
  \par\smallskip\noindent\emph{Definition.} Spanning subgraph with edge set \(E(H_{\mathcal P})=\bigcup_{B\in\mathcal P}\binom{B}{2}\).
  \item \texttt{\textbackslash{}\allowbreak{}xi\_\allowbreak{}\{B,\allowbreak{}k\}\allowbreak{}} --- \texttt{block\_\allowbreak{}layer\_\allowbreak{}rademacher\_\allowbreak{}sign}; \(\{-1,1\}\); \texttt{least\_\allowbreak{}favourable\_\allowbreak{}prior\_\allowbreak{}coordinate}
  \par\smallskip\noindent\emph{Definition.} Independent Rademacher sign for clique-partition block \(B\) and target layer \(k\).
  \item \texttt{\textbackslash{}\allowbreak{}alpha(G)} --- \texttt{independence\_\allowbreak{}number}; \(\{1,\ldots,n\}\); \texttt{published\_\allowbreak{}graph\_\allowbreak{}functional}
  \par\smallskip\noindent\emph{Definition.} Maximum cardinality of an independent vertex set in \(G\).
  \item \texttt{d\_\allowbreak{}\{\textbackslash{}\allowbreak{}mathrm\{avg\}\allowbreak{}\}\allowbreak{}(G)} --- \texttt{average\_\allowbreak{}degree}; \([0,n-1]\); \texttt{published\_\allowbreak{}graph\_\allowbreak{}functional}
  \par\smallskip\noindent\emph{Definition.} Average vertex degree of the simple graph \(G\).
  \item \texttt{\textbackslash{}\allowbreak{}lambda(G)} --- \texttt{adjacency\_\allowbreak{}spectral\_\allowbreak{}radius}; \([0,n-1]\); \texttt{published\_\allowbreak{}graph\_\allowbreak{}functional}
  \par\smallskip\noindent\emph{Definition.} Largest eigenvalue of the adjacency matrix of the simple graph \(G\).
  \item \texttt{\textbackslash{}\allowbreak{}rho} --- \texttt{fixed\_\allowbreak{}dense\_\allowbreak{}clique\_\allowbreak{}fraction}; \((0,1)\); \texttt{rate\_\allowbreak{}family\_\allowbreak{}parameter}
  \par\smallskip\noindent\emph{Definition.} \texttt{Fixed limiting fraction used in the dense-\allowbreak{}clique-\allowbreak{}plus-\allowbreak{}isolates calibration family.}
  \item \texttt{\textbackslash{}\allowbreak{}omega(G)} --- \texttt{clique\_\allowbreak{}number}; \(\{1,\ldots,n\}\); \texttt{witnessed\_\allowbreak{}graph\_\allowbreak{}functional}
  \par\smallskip\noindent\emph{Definition.} Maximum cardinality of a vertex subset inducing a complete subgraph of \(G\).
  \item \texttt{\textbackslash{}\allowbreak{}gamma} --- \texttt{polynomial\_\allowbreak{}clique\_\allowbreak{}size\_\allowbreak{}exponent}; \((0,1]\); \texttt{rate\_\allowbreak{}family\_\allowbreak{}parameter}
  \par\smallskip\noindent\emph{Definition.} Exponent in the calibration family with clique size \(m_n=\lfloor\rho n^{\gamma}\rfloor\).
  \item \texttt{\textbackslash{}\allowbreak{}beta\_\allowbreak{}\{\textbackslash{}\allowbreak{}mathrm\{DTE\}\allowbreak{}\}\allowbreak{}(G)} --- \texttt{cross\_\allowbreak{}layer\_\allowbreak{}balanced\_\allowbreak{}biclique\_\allowbreak{}parameter}; \(\{1,\ldots,n\}\); \texttt{arbitrary\_\allowbreak{}graph\_\allowbreak{}lower\_\allowbreak{}certificate}
  \par\smallskip\noindent\emph{Definition.} \texttt{Maximum balanced size of a nonempty ordered control-\allowbreak{}support and treated-\allowbreak{}support pair satisfying the common closed-\allowbreak{}neighborhood containment.}
\end{itemize}

\section{Assumptions}

\begin{assumption}[ass:arbitrary-neighborhood-interference]\label{ass:arbitrary-neighborhood-interference}
For every \(i\in V_n\) and \(z,z^{\prime}\in\mathcal Z_n\), equality \(z|_{\widetilde N_i}=z^{\prime}|_{\widetilde N_i}\) implies \(Y_i(z)=Y_i(z^{\prime})\).
\par\smallskip\noindent\emph{Standard} (Arbitrary neighborhood interference, \cite{KandirosHarshawSavje2026}).
\end{assumption}

\begin{assumption}[ass:control-layer-second-moment]\label{ass:control-layer-second-moment}
The control-exposure outcomes satisfy \(n^{-1}\sum_{i\in V_n}u_{i,0}^2\le1\).
\par\smallskip\noindent\emph{Standard} (Published \(q=2\) control-layer moment restriction, \cite{KandirosHarshawSavje2026}).
\end{assumption}

\begin{assumption}[ass:treated-layer-second-moment]\label{ass:treated-layer-second-moment}
The isolated-treatment outcomes satisfy \(n^{-1}\sum_{i\in V_n}u_{i,1}^2\le1\).
\par\smallskip\noindent\emph{Standard} (Published \(q=2\) treated-layer moment restriction, \cite{KandirosHarshawSavje2026}).
\end{assumption}

\section{Definitions}

\begin{definition}[def:block-clique-family]\label{def:block-clique-family}
\par\smallskip\noindent\emph{Defined object.} \texttt{Equal block-\allowbreak{}clique family}
\par\smallskip\noindent\emph{Construction.} For \(b\ge1\), \(s\ge2\), and \(n=bs\), fix \(\mathcal B_{b,s}=\{B_1,\ldots,B_b\}\) and set \(E(G_{b,s})=\{\{i,i^{\prime}\}:i\ne i^{\prime}\text{ and }i,i^{\prime}\in B_r\text{ for some }r\}\).
\par\smallskip\noindent\emph{Inputs.} \texttt{b}; \texttt{s}; \texttt{\textbackslash{}\allowbreak{}mathcal B\_\allowbreak{}\{b,\allowbreak{}s\}\allowbreak{}}.
\end{definition}

\begin{definition}[def:published-q2-class]\label{def:published-q2-class}
\par\smallskip\noindent\emph{Defined object.} \texttt{Published second-\allowbreak{}moment ANI class}
\par\smallskip\noindent\emph{Construction.} \(\mathcal M_2(G_{b,s})=\{Y=(Y_i(z))_{i\in V_n,z\in\mathcal Z_n}:Y\text{ satisfies }\text{ass:arbitrary-neighborhood-interference},\text{ ass:control-layer-second-moment},\text{ and ass:treated-layer-second-moment}\}\).
\par\smallskip\noindent\emph{Member properties.} \texttt{ass:\allowbreak{}arbitrary-\allowbreak{}neighborhood-\allowbreak{}interference}; \texttt{ass:\allowbreak{}control-\allowbreak{}layer-\allowbreak{}second-\allowbreak{}moment}; \texttt{ass:\allowbreak{}treated-\allowbreak{}layer-\allowbreak{}second-\allowbreak{}moment}.
\end{definition}

\begin{definition}[def:layerwise-parameter-set]\label{def:layerwise-parameter-set}
\par\smallskip\noindent\emph{Defined object.} \texttt{Separate-\allowbreak{}layer Euclidean parameter set}
\par\smallskip\noindent\emph{Construction.} \[\Theta_{b,s}=\left\{u\in\mathbb R^{V_n\times\{0,1\}}:n^{-1}\sum_{i\in V_n}u_{i,0}^2\le1\text{ and }n^{-1}\sum_{i\in V_n}u_{i,1}^2\le1\right\}.\]
\par\smallskip\noindent\emph{Inputs.} \texttt{b}; \texttt{s}; \texttt{n}; \texttt{V\_\allowbreak{}n}.
\end{definition}

\begin{definition}[def:unrestricted-minimax-risk]\label{def:unrestricted-minimax-risk}
\par\smallskip\noindent\emph{Defined object.} \texttt{Unrestricted block-\allowbreak{}clique DTE minimax risk}
\par\smallskip\noindent\emph{Construction.} \[R_2(G_{b,s})=\inf_{P\in\mathcal D_n,\,\widehat\tau\in\Gamma_n}\sup_{Y\in\mathcal M_2(G_{b,s})}\mathbb E_{z\sim P}\!\left[\left\{\widehat\tau(z,Y(z))-\tau_{\mathrm{DTE}}(Y)\right\}^2\right].\]
\par\smallskip\noindent\emph{Inputs.} \texttt{G\_\allowbreak{}\{b,\allowbreak{}s\}\allowbreak{}}.
\end{definition}

\begin{definition}[def:block-clique-conflict-orbits]\label{def:block-clique-conflict-orbits}
\par\smallskip\noindent\emph{Defined object.} \texttt{Block-\allowbreak{}clique maximal-\allowbreak{}exposure orbits}
\par\smallskip\noindent\emph{Construction.} For \(c=0,\ldots,b\), \(\mathcal S_c\) contains the maximal exposure sets obtained by choosing \(C\subseteq\{1,\ldots,b\}\) with \(|C|=c\), revealing every \((i,0)\) in blocks \(r\in C\), and revealing one \((J_r,1)\) in every block \(r\notin C\).
\par\smallskip\noindent\emph{Inputs.} \texttt{b}; \texttt{s}; \texttt{H\_\allowbreak{}\{b,\allowbreak{}s\}\allowbreak{}}.
\end{definition}

\begin{definition}[def:invariant-procedure]\label{def:invariant-procedure}
\par\smallskip\noindent\emph{Defined object.} \texttt{Invariant block-\allowbreak{}clique design and estimator}
\par\smallskip\noindent\emph{Construction.} Draw \(c\sim p\); conditional on \(c\), draw \(C\) uniformly among the \(c\)-subsets of \(\{1,\ldots,b\}\), set all units in blocks \(r\in C\) to control, and independently draw \(J_r\) uniformly in each block \(r\notin C\) and treat only that unit. On the realized assignment set \[\delta_{p,A,D}=A_c\sum_{r\notin C}u_{J_r,1}+D_c\sum_{r\in C}\sum_{i\in B_r}u_{i,0}.\]
\par\smallskip\noindent\emph{Inputs.} \texttt{p}; \texttt{A\_\allowbreak{}c}; \texttt{D\_\allowbreak{}c}.
\end{definition}

\begin{definition}[def:mode-risk-matrices]\label{def:mode-risk-matrices}
\par\smallskip\noindent\emph{Defined object.} \texttt{Constant,\allowbreak{} between-\allowbreak{}block,\allowbreak{} and within-\allowbreak{}block risk matrices}
\par\smallskip\noindent\emph{Construction.} Write \(W_c=(w^c_{\alpha\beta})_{\alpha,\beta=0}^2\), set \(\rho_c=c(b-c)/\{b(b-1)\}\) for \(b>1\) and \(\rho_c=0\) for \(b=1\), and define \[K_0(W)=\sum_{c=0}^b\begin{pmatrix}w^c_{00}+2scw^c_{02}+s^2c^2w^c_{22}&-w^c_{00}+(b-c)w^c_{01}-scw^c_{02}+sc(b-c)w^c_{12}\\-w^c_{00}+(b-c)w^c_{01}-scw^c_{02}+sc(b-c)w^c_{12}&w^c_{00}-2(b-c)w^c_{01}+(b-c)^2w^c_{11}\end{pmatrix},\] \[K_{\mathrm B}(W)=\sum_{c=0}^b b\rho_c\begin{pmatrix}s^2w^c_{22}&-sw^c_{12}\\-sw^c_{12}&w^c_{11}\end{pmatrix},\qquad k_{\mathrm W}(W)=\sum_{c=0}^b(b-c)w^c_{11}.\]
\par\smallskip\noindent\emph{Inputs.} \texttt{b}; \texttt{s}; \texttt{W\_\allowbreak{}c}.
\end{definition}

\begin{definition}[def:mode-sdp-value]\label{def:mode-sdp-value}
\par\smallskip\noindent\emph{Defined object.} \texttt{Explicit finite block-\allowbreak{}clique mode SDP}
\par\smallskip\noindent\emph{Construction.} \[V_{b,s}=\min_{(W_c)_{c=0}^b,\lambda_0,\lambda_1}\lambda_0+\lambda_1\] subject to \(W_c\succeq0\), \(\sum_{c=0}^bw^c_{00}=1\), \(\lambda_0,\lambda_1\ge0\), and \[\operatorname{diag}(\lambda_0,\lambda_1)-K_0(W)\succeq0,\qquad\operatorname{diag}(\lambda_0,\lambda_1)-K_{\mathrm B}(W)\succeq0,\qquad\lambda_1\ge k_{\mathrm W}(W).\]
\par\smallskip\noindent\emph{Inputs.} \texttt{b}; \texttt{s}.
\end{definition}

\begin{definition}[def:posterior-collision-lift]\label{def:posterior-collision-lift}
\par\smallskip\noindent\emph{Defined object.} \texttt{Finite orbit posterior-\allowbreak{}collision lift}
\par\smallskip\noindent\emph{Construction.} Fix a primal--dual optimum of \(\text{def:mode-sdp-value}\). For each nonzero dual constant-mode and between-block matrix, take the spectral factor whose eigenvalues are in decreasing order and whose eigenvectors use the lexicographic sign convention in the displayed two-dimensional mode basis; take the nonnegative square root for the within-block scalar. Embed the resulting at most five columns using \(g_0\), \(g_{\mathrm B}\), and \(g_{\mathrm W}\). Divide each embedded vector \(v\) by \(\max\{1,\|v_0\|_2/\sqrt n,\|v_1\|_2/\sqrt n\}\). For every \(c\in\mathcal C^{\star}\), form \(P_cv\pm(I-P_c)v\), close these atoms under \(\mathfrak G_{b,s}\), and discard duplicates; if every dual column is zero, use the singleton dictionary containing the zero vector. This finite set is \(\mathcal A_{b,s}^{\star}\). On its simplex of nonnegative weights that sum to one and are constant on \(\mathfrak G_{b,s}\)-orbits, choose the lexicographically first maximizer of \(\min_{0\le c\le b}B_c(\sum_hq_h\delta_{u^{(h)}})\). Denote the resulting prior by \(\Pi_{b,s}^{\star}=\sum_hq_h\delta_{u^{(h)}}\). The construction returns this finite, directly checkable maximin posterior certificate; it does not assert that its minimum posterior risk equals \(V_{b,s}\).
\par\smallskip\noindent\emph{Inputs.} \texttt{b}; \texttt{s}; \texttt{V\_\allowbreak{}\{b,\allowbreak{}s\}\allowbreak{}}.
\end{definition}

\begin{definition}[def:finite-posterior-risk]\label{def:finite-posterior-risk}
\par\smallskip\noindent\emph{Defined object.} \texttt{Orbitwise finite-\allowbreak{}prior posterior risk}
\par\smallskip\noindent\emph{Construction.} For a finite prior \(\Pi=\sum_hq_h\delta_{u^{(h)}}\) and the representative \(S_c\in\mathcal S_c\), partition its atoms into cells \(\mathcal Q\) with identical revealed subvectors \(u^{(h)}_{S_c}\), set \(m_{\mathcal Q}=(\sum_{h\in\mathcal Q}q_h)^{-1}\sum_{h\in\mathcal Q}q_h\,n^{-1}\sum_{i\in V_n}(u^{(h)}_{i,1}-u^{(h)}_{i,0})\), and define \[B_c(\Pi)=\sum_{\mathcal Q}\sum_{h\in\mathcal Q}q_h\left\{n^{-1}\sum_{i\in V_n}(u^{(h)}_{i,1}-u^{(h)}_{i,0})-m_{\mathcal Q}\right\}^2.\]
\par\smallskip\noindent\emph{Inputs.} \texttt{\textbackslash{}\allowbreak{}Pi\_\allowbreak{}\{b,\allowbreak{}s\}\allowbreak{}\textasciicircum{}\{\textbackslash{}\allowbreak{}star\}\allowbreak{}}; \texttt{\textbackslash{}\allowbreak{}mathcal S\_\allowbreak{}c}.
\end{definition}

\begin{definition}[def:general-graph-certificate-handle]\label{def:general-graph-certificate-handle}
\par\smallskip\noindent\emph{Defined object.} \texttt{Orbit-\allowbreak{}aware saddle column-\allowbreak{}generation handle}
\par\smallskip\noindent\emph{Construction.} \texttt{Use posterior-\allowbreak{}collision orbit atoms as lower columns,\allowbreak{} alternate them with maximal-\allowbreak{}compatible-\allowbreak{}exposure separation on the design side,\allowbreak{} and apply Carathéodory reduction to the two layer moment map while retaining exact posterior-\allowbreak{}cell constraints.}
\par\smallskip\noindent\emph{Inputs.} \texttt{H\_\allowbreak{}\{\textbackslash{}\allowbreak{}mathrm\{DTE\}\allowbreak{}\}\allowbreak{}(G)}; \texttt{\textbackslash{}\allowbreak{}Theta\_\allowbreak{}2(G)}.
\end{definition}

\begin{definition}[def:general-dte-primitives]\label{def:general-dte-primitives}
\par\smallskip\noindent\emph{Defined object.} \texttt{General-\allowbreak{}graph DTE primitives}
\par\smallskip\noindent\emph{Construction.} For a simple graph \(G\) on \(V_n\), define the closed neighborhood
\[
\widetilde N_i(G)=\{i\}\cup\{j\in V_n:\{i,j\}\in E(G)\}.
\]
Let \(\mathbf 0\in\mathcal Z_n\) be the all-control assignment and let \(\mathbf e_i\in\mathcal Z_n\) treat only unit \(i\). Define
\[
\tau_{\mathrm{DTE}}^G(Y)
=\frac1n\sum_{i\in V_n}\{Y_i(\mathbf e_i)-Y_i(\mathbf 0)\}.
\]
The superscript records the model relative to which the schedule is restricted; the displayed finite-population contrast uses the same two physical assignments for every graph.
\par\smallskip\noindent\emph{Inputs.} \texttt{G}; \texttt{V\_\allowbreak{}n}; \texttt{n}.
\end{definition}

\begin{definition}[def:general-published-q2-class]\label{def:general-published-q2-class}
\par\smallskip\noindent\emph{Defined object.} \texttt{General published second-\allowbreak{}moment ANI class}
\par\smallskip\noindent\emph{Construction.} For a simple graph \(G\) on \(V_n\), define
\[
\mathcal M_2(G)=\left\{Y:\begin{array}{l}
\text{for every }i\in V_n\text{ and }z,z'\in\mathcal Z_n,
\ z|_{\widetilde N_i(G)}=z'|_{\widetilde N_i(G)}\Rightarrow Y_i(z)=Y_i(z'),\\[2mm]
\displaystyle \frac1n\sum_{i\in V_n}Y_i(\mathbf 0)^2\leq1,
\qquad
\displaystyle \frac1n\sum_{i\in V_n}Y_i(\mathbf e_i)^2\leq1
\end{array}\right\}.
\]
Thus \(\mathcal M_2(G)\) is exactly the published arbitrary-neighborhood-interference class with separate \(q=2\) restrictions on the all-control and isolated-treatment target layers.
\par\smallskip\noindent\emph{Inputs.} \texttt{G}.
\end{definition}

\begin{definition}[def:general-dte-minimax-risk]\label{def:general-dte-minimax-risk}
\par\smallskip\noindent\emph{Defined object.} \texttt{General unrestricted DTE minimax risk}
\par\smallskip\noindent\emph{Construction.} For a simple graph \(G\) on \(V_n\), define the unrestricted design-based DTE minimax risk
\[
R_2(G)=\inf_{P\in\mathcal D_n,\,\widehat\tau\in\Gamma_n}
\sup_{Y\in\mathcal M_2(G)}
\mathbb E_{z\sim P}\!\left[
\left\{\widehat\tau(z,Y(z))-\tau_{\mathrm{DTE}}^G(Y)\right\}^2
\right].
\]
The only randomness is \(z\sim P\); the graph, population, and schedule are fixed.
\par\smallskip\noindent\emph{Inputs.} \texttt{G}.
\end{definition}

\begin{definition}[def:clique-partition-subgraph]\label{def:clique-partition-subgraph}
\par\smallskip\noindent\emph{Defined object.} \texttt{Clique partition and spanning clique subgraph}
\par\smallskip\noindent\emph{Construction.} A clique partition \(\mathcal P\) of \(G\) is a partition of \(V_n\) into nonempty blocks \(B\) such that every block induces a clique of \(G\). Its spanning clique subgraph is
\[
V(H_{\mathcal P})=V_n,
\qquad
E(H_{\mathcal P})=\bigcup_{B\in\mathcal P}\binom{B}{2}.
\]
In particular, \(H_{\mathcal P}\subseteq G\), and singleton blocks are allowed.
\par\smallskip\noindent\emph{Inputs.} \texttt{G}; \texttt{\textbackslash{}\allowbreak{}mathcal P}.
\end{definition}

\begin{definition}[def:balanced-biclique-dte-parameter]\label{def:balanced-biclique-dte-parameter}
\par\smallskip\noindent\emph{Defined object.} \texttt{Cross-\allowbreak{}layer balanced-\allowbreak{}biclique DTE parameter}
\par\smallskip\noindent\emph{Construction.} For every simple graph \(G\) on the nonempty finite population \(V_n\), define
\[
\beta_{\mathrm{DTE}}(G)
=
\max_{\substack{\varnothing\ne A\subseteq V_n,\ \varnothing\ne B\subseteq V_n\\
B\subseteq\bigcap_{i\in A}\widetilde N_i(G)}}
\min\{|A|,|B|\}.
\]
The maximization is over ordered cross-layer support pairs: \(A\) supports the all-control layer and \(B\) supports the isolated-treatment layer.
\par\smallskip\noindent\emph{Inputs.} \texttt{G}.
\end{definition}

\section{Main results}

\begin{lemma}[lem:published-conflict-calibration]\label{lem:published-conflict-calibration}
The published maximal-independent-set identity specializes to \(G_{b,s}\): every maximal compatible exposure set is realized by an assignment, and the unrestricted risk depends only on the two target-outcome layers and the family \(\bigcup_{c=0}^b\mathcal S_c\).
\end{lemma}
\begin{proof}
\textbf{Definitions.} For block \(B_r\), write
\[
\mathsf C_r=\{(i,0):i\in B_r\},
\qquad
\mathsf T_r=\{(i,1):i\in B_r\}.
\]
The exposure \((i,0)\) requires all coordinates in \(B_r\) to be zero, whereas \((i,1)\) requires coordinate \(i\) to be one and every other coordinate in \(B_r\) to be zero. These descriptions follow from the two target-exposure definitions and \(\text{ass:arbitrary-neighborhood-interference}\), because the closed neighborhood of every unit in \(B_r\) is exactly \(B_r\).

\textbf{Compatibility inside one block.} All vertices in \(\mathsf C_r\) are jointly realized by the all-control assignment on \(B_r\). Distinct vertices of \(\mathsf T_r\) are incompatible: realizing \((i,1)\) sets coordinate \(i\) to one and every other coordinate of \(B_r\) to zero, whereas realizing \((j,1)\), for \(j\ne i\), requires the reverse at coordinates \(i,j\). Every member of \(\mathsf T_r\) is also incompatible with every member of \(\mathsf C_r\), since the former requires one treated coordinate and the latter requires none. Exposures belonging to distinct blocks are compatible because the blocks are disjoint. Consequently, a maximal compatible set chooses, independently in each block, either all of \(\mathsf C_r\) or one member of \(\mathsf T_r\). This is precisely a member of one of the families \(\mathcal S_c\).

\textbf{Assignment realization.} Given such a maximal set, let \(C\) be its all-control blocks. Assign every unit in a block in \(C\) to control. In each block outside \(C\), assign only its selected label \(J_r\) to treatment. The preceding compatibility calculation shows that the revealed target coordinates are exactly the prescribed maximal set. Thus every maximal compatible exposure set is realized by a lawful binary assignment.

\textbf{Reduction of the minimax game.} For an assignment \(z\), let \(R_z\) be the set of target coordinates revealed by \(z\). The preceding argument shows that \(R_z\) is compatible. Choose a maximal extension \(M(R_z)\supseteq R_z\), and realize it by the blockwise assignment just constructed. Replacing a design by its pushforward through \(z\mapsto M(R_z)\) cannot increase its optimal risk: conditional on a maximal extension, average the old estimator over its preimage and apply Jensen's inequality to squared loss. The new observation contains every target coordinate used by the old estimator.

It remains to justify discarding non-target outcomes. For the lower inequality, restrict the supremum to schedules for which every non-target exposure outcome is zero. Such schedules satisfy arbitrary neighborhood interference and are indexed by \(u\in\Theta_{b,s}\); on this subclass the data under \(z\) consist only of \(R_z\), the revealed coordinates \(u_{R_z}\), and fixed zeros. Hence every network procedure induces a procedure in the target-coordinate game. For the reverse inequality, realize each maximal compatible set by the displayed assignment and use an estimator that ignores every non-target outcome. The two induced risks agree pointwise in \(u\). Taking first the supremum over the two layer balls and then the infimum over designs and measurable estimators proves that the unrestricted risk depends only on \(u\) and the family \(\bigcup_{c=0}^b\mathcal S_c\).

This is the block-clique specialization of the maximal-independent-set reduction in Kandiros, Harshaw, and S\"avje (2026, Proposition~2.2 and Appendix~A.5); all specialization steps used here were derived directly.
\end{proof}
\par\smallskip\noindent \textit{Justification.} This calibration identifies the lawful target-coordinate game and proves that maximal compatible exposure sets are realizable on equal block cliques. \textit{Gap.} KandirosHarshawSavje2026 proves the general conflict-graph identity; the node records only its direct equal-block-clique specialization. \textit{Consumer.} Both upper realizability and the arbitrary-assignment Bayes lower arguments use this reduction without treating it as the new contribution.

\begin{lemma}[lem:three-mode-risk-sdp]\label{lem:three-mode-risk-sdp}
For every invariant procedure \(\delta_{p,A,D}\), its risk quadratic form decomposes orthogonally into the constant block-mean mode \(K_0(W)\), \(b-1\) identical between-block zero-sum modes \(K_{\mathrm B}(W)\), and \(b(s-1)\) treated within-block contrast modes with eigenvalue \(k_{\mathrm W}(W)\); control within-block contrasts have eigenvalue zero. Maximization over the two separate unit Euclidean balls is exactly dual to the two multipliers \(\lambda_0,\lambda_1\). Every optimum of the convex moment relaxation admits rank-one purification \(W_c=p_c(1,A_c,D_c)^{\mathsf T}(1,A_c,D_c)\) without increasing risk.
\end{lemma}
\begin{proof}
\textbf{Error coefficients.} Fix \(c\), \(C\), and the sampled labels \((J_r:r\notin C)\). Write \(u_k=(u_{i,k})_{i\in V_n}\) and put
\[
a_{0,c,C}=n^{-1}\mathbf 1_{V_n}+D_c\sum_{r\in C}\mathbf 1_{B_r},
\qquad
a_{1,c,C,J}=-n^{-1}\mathbf 1_{V_n}+A_c\sum_{r\notin C}\mathbf e_{J_r}.
\]
The estimation error is exactly
\[
\delta_{p,A,D}-\tau_{\mathrm{DTE}}(Y)
=\langle a_{0,c,C},u_0\rangle+\langle a_{1,c,C,J},u_1\rangle.
\]
Thus its design risk is a positive-semidefinite quadratic form in \((u_0,u_1)\).

\textbf{Orthogonal modes.} Decompose \(\mathbb R^{V_n}\) into the orthogonal sum of the constant vector, the \((b-1)\)-dimensional space of vectors that are constant on each block and whose block constants sum to zero, and the \(b(s-1)\)-dimensional space of within-block zero-sum vectors. Conditional uniformity of \(C\) and the labels \(J_r\) makes the risk operator equivariant under block and within-block permutations. Directly averaging the displayed error coefficients therefore makes the three spaces invariant, makes all between-block copies identical, and makes all treated within-block copies identical. The control coefficient is constant inside every block, so its inner product with each control within-block contrast is zero.

Scale \(u_k=\sqrt n\,x_k\), so the two restrictions defining \(\Theta_{b,s}\) become \(\|x_0\|_2\le1\) and \(\|x_1\|_2\le1\). On the normalized constant vector \(g_0\), the two scaled error coefficients are
\[
1+scD_c,
\qquad
-1+(b-c)A_c.
\]
If
\[
L_{0c}=\begin{pmatrix}1&0&sc\\-1&b-c&0\end{pmatrix},
\]
then the constant-mode contribution is \(L_{0c}W_cL_{0c}^{\mathsf T}\). Summing over \(c\) gives exactly \(K_0(W)\).

For \(b>1\), use the normalized contrast between two distinct blocks. Conditional on \(c\),
\[
\Pr(1\in C,\,2\notin C)=\Pr(2\in C,\,1\notin C)
=\rho_c=\frac{c(b-c)}{b(b-1)}.
\]
When exactly one of the two blocks belongs to \(C\), the scaled two-layer coefficient vector is, up to orientation, proportional to \((sD_c,-A_c)\); the normalization \(n=bs\) and the two orientations give the factor \(b\rho_c\). Therefore the common between-block matrix is
\[
\sum_{c=0}^b b\rho_c
\begin{pmatrix}
s^2p_cD_c^2&-sp_cA_cD_c\\
-sp_cA_cD_c&p_cA_c^2
\end{pmatrix}
=K_{\mathrm B}(W).
\]
For \(b=1\) the between-block subspace is zero-dimensional and the definition sets \(\rho_c=0\), so the same conclusion holds. For a normalized treated within-block contrast, uniformity of \(J_r\) gives mean squared sampled coordinate \(1/s\), and its block is outside \(C\) with probability \((b-c)/b\). Multiplication by \(n=bs\) gives
\[
k_{\mathrm W}(W)=\sum_{c=0}^b(b-c)p_cA_c^2.
\]
The same conditional averages show that cross terms between distinct orthogonal modes vanish. These identities prove the asserted decomposition for rank-one moment matrices.

\textbf{Exact maximization over the product of balls.} Let \(Q(W)\) be the block-diagonal matrix with one copy of \(K_0(W)\), \(b-1\) copies of \(K_{\mathrm B}(W)\), the treated within-block scalars \(k_{\mathrm W}(W)\), and zero control within-block scalars. The worst-case risk is
\[
\sup_{\|x_0\|_2\le1,\,\|x_1\|_2\le1}x^{\mathsf T}Q(W)x.
\]
Lift \(xx^{\mathsf T}\) to \(X\succeq0\) subject to the two layer-trace constraints \(\operatorname{tr}(E_kX)\le1\), \(k=0,1\). This relaxation is exact. Indeed, if an extreme feasible \(X\) had rank \(r\ge2\), symmetric perturbations supported on its range would have dimension \(r(r+1)/2\ge3\). At most two layer-trace inequalities are active, so a nonzero such perturbation can annihilate every active trace functional. Perturbations of both signs with sufficiently small magnitude preserve positive semidefiniteness and every inactive inequality, contradicting extremality. Hence every nonzero extreme point has rank one, and compactness supplies a rank-one maximizer.

The dual is strictly feasible by choosing both multipliers larger than \(\|Q(W)\|_{\mathrm{op}}\). Finite-dimensional semidefinite strong duality yields
\[
\sup x^{\mathsf T}Q(W)x
=\min_{\lambda_0,\lambda_1\ge0}
\left\{\lambda_0+\lambda_1:
\lambda_0E_0+\lambda_1E_1-Q(W)\succeq0\right\}.
\]
Block diagonality makes the final constraint equivalent to
\[
\operatorname{diag}(\lambda_0,\lambda_1)-K_0(W)\succeq0,
\qquad
\operatorname{diag}(\lambda_0,\lambda_1)-K_{\mathrm B}(W)\succeq0,
\qquad
\lambda_1\ge k_{\mathrm W}(W).
\]

\textbf{Purification of relaxed moments.} For a relaxed block \(W_c\succeq0\), let \(p_c=w^c_{00}\). If \(p_c>0\), define
\[
\bar A_c=\frac{w^c_{01}}{p_c},
\qquad
\bar D_c=\frac{w^c_{02}}{p_c},
\qquad
\bar W_c=p_c(1,\bar A_c,\bar D_c)^{\mathsf T}(1,\bar A_c,\bar D_c).
\]
The Schur complement gives \(W_c-\bar W_c\succeq0\). If \(p_c=0\), positive semidefiniteness forces the first row and column of \(W_c\) to vanish; set \(\bar W_c=0\), so again \(W_c-\bar W_c\succeq0\). The identities
\[
K_0(W)=\sum_cL_{0c}W_cL_{0c}^{\mathsf T},
\qquad
K_{\mathrm B}(W)=\sum_cb\rho_cL_{\mathrm B}W_cL_{\mathrm B}^{\mathsf T},
\qquad
k_{\mathrm W}(W)=\sum_c(b-c)e_1^{\mathsf T}W_ce_1,
\]
where
\[
L_{\mathrm B}=\begin{pmatrix}0&0&s\\0&-1&0\end{pmatrix}
\]
and \(e_1\) selects the \(A_c\)-coordinate, show by congruence monotonicity that
\[
K_0(\bar W)\preceq K_0(W),
\qquad
K_{\mathrm B}(\bar W)\preceq K_{\mathrm B}(W),
\qquad
k_{\mathrm W}(\bar W)\le k_{\mathrm W}(W).
\]
The constraint \(\sum_cp_c=1\) is unchanged. Thus every feasible relaxed point has a rank-one purification with no larger epigraph value. Conversely every \((p,A,D)\) produces the displayed rank-one moments. The SDP therefore equals the infimum of the worst-case risks over the invariant linear class, and every SDP optimum can be purified as asserted.
\end{proof}
\par\smallskip\noindent \textit{Justification.} This lemma gives the exact finite optimization and rank-one extraction for the invariant linear design--estimator class. \textit{Gap.} Published graph-summary rates do not diagonalize this block-clique linear risk, and the lemma does not identify nonlinear invariant estimators with the linear subclass. \textit{Consumer.} The main theorem uses it to output a realizable randomization, explicit estimator coefficients, and a computable unrestricted upper-risk certificate.

\begin{theorem}[thm:block-clique-exact-saddle]\label{thm:block-clique-exact-saddle}
For every \(b\ge1\) and \(s\ge2\), the explicit finite mode SDP is exactly the minimax value within the invariant linear class:
\[
V_{b,s}=
\inf_{p,A,D}\sup_{u\in\Theta_{b,s}}
\mathbb E\!\left[\left\{\delta_{p,A,D}-\tau_{\mathrm{DTE}}(Y)\right\}^2\right],
\]
where \(p\) ranges over the orbit simplex and \(A_c,D_c\in\mathbb R\). If \(W_c^{\star}=p_c^{\star}(1,A_c^{\star},D_c^{\star})^{\mathsf T}(1,A_c^{\star},D_c^{\star})\) is a purified optimum, then \(\delta_{p^{\star},A^{\star},D^{\star}}\) is a realizable assignment law and measurable estimator, and therefore \(R_2(G_{b,s})\le V_{b,s}\). The corrected posterior-collision search returns an atomwise-feasible finite prior \(\Pi_{b,s}^{\star}\), supported on explicit \(\mathfrak G_{b,s}\)-orbits with \(L_{b,s}\le10(b+1)|\mathfrak G_{b,s}|\), and supplies the computable sandwich
\[
\min_{0\le c\le b}B_c(\Pi_{b,s}^{\star})
\le R_2(G_{b,s})\le V_{b,s}.
\]
If the finite posterior-cell calculation verifies \(\min_cB_c(\Pi_{b,s}^{\star})=V_{b,s}\), then the displayed procedure and prior form an exact all-estimator saddle. Without that verification, no equality between \(R_2(G_{b,s})\) and \(V_{b,s}\) is asserted.
\end{theorem}
\begin{proof}
\textbf{Proof.}

\textbf{Invariant-linear value and realizability.}
By \(\mathrm{lem{:}three\text{-}mode\text{-}risk\text{-}sdp}\), the worst-case risk of every procedure \(\delta_{p,A,D}\) is the value of the two-multiplier epigraph at its rank-one moment matrices. Conversely, the purification part of that lemma converts every feasible SDP point into rank-one conditional moments without increasing the constant, between-block, or within-block mode risks. Therefore
\[
V_{b,s}
=\inf_{p,A,D}\sup_{u\in\Theta_{b,s}}
\mathbb E\!\left[\left\{\delta_{p,A,D}-\tau_{\mathrm{DTE}}(Y)\right\}^2\right].
\tag{1}
\]
The randomization defining \(\delta_{p,A,D}\) first draws \(c\), then a uniformly random \(c\)-subset \(C\), and finally one uniformly random label in every block outside \(C\). It assigns every block in \(C\) to all control and treats only the sampled label in every other block. By \(\mathrm{lem{:}published\text{-}conflict\text{-}calibration}\), each resulting binary assignment lawfully realizes its stated maximal target-exposure pattern. The displayed linear rule is measurable in the realized assignment and observed outcomes. Hence a purified optimizer is an admissible member of \(\mathcal D_n\times\Gamma_n\). Since the unrestricted infimum ranges over this invariant linear subclass,
\[
R_2(G_{b,s})\leq V_{b,s}.
\tag{2}
\]

\textbf{Feasibility and size of the posterior-collision certificate.}
The normalization in \(\mathrm{def{:}posterior\text{-}collision\text{-}lift}\) puts each embedded vector \(v\) in the two separate layer balls. The coordinate projection \(P_c\) preserves the layer decomposition and is orthogonal. Thus, for each \(k\in\{0,1\}\),
\[
\left\|P_cv_k\mathbin{\pm}(I-P_c)v_k\right\|_2^2
=\|P_cv_k\|_2^2+\|(I-P_c)v_k\|_2^2
=\|v_k\|_2^2\leq n.
\tag{3}
\]
Every collision atom therefore lies in \(\Theta_{b,s}\). Block permutations, within-block permutations, and simultaneous sign reversal preserve both layer norms, so closure under \(\mathfrak G_{b,s}\) preserves feasibility.

A positive-semidefinite constant-mode matrix has at most two spectral columns, a positive-semidefinite between-block matrix has at most two, and the nonnegative within-block scalar has at most one. Hence there are at most five columns before collision. The two collision signs and at most \(b+1\) active orbit indices produce at most \(10(b+1)\) atoms before group closure, and group closure multiplies cardinality by at most \(|\mathfrak G_{b,s}|\). Discarding duplicates can only reduce this number, so
\[
L_{b,s}\leq10(b+1)|\mathfrak G_{b,s}|.
\tag{4}
\]
If all dual columns vanish, the construction instead uses the feasible zero atom and the same bound remains valid.

The orbit-constant weight vectors form a nonempty compact simplex. For a posterior cell \(\mathcal Q\), put
\[
\tau_h=\frac1n\sum_{i\in V_n}\bigl(u^{(h)}_{i,1}-u^{(h)}_{i,0}\bigr).
\]
When the cell mass is positive, its contribution to posterior risk is
\[
\sum_{h\in\mathcal Q}q_h\tau_h^2
-\frac{\left(\sum_{h\in\mathcal Q}q_h\tau_h\right)^2}
{\sum_{h\in\mathcal Q}q_h}.
\tag{5}
\]
Assign value zero when the cell mass is zero. Cauchy--Schwarz bounds the absolute value of the quotient by a fixed constant times the cell mass, so (5) extends continuously to zero mass. Thus every \(B_c\), their finite minimum, and the lexicographically selected maximizer defining \(\Pi_{b,s}^{\star}\) are well defined and directly computable by finite posterior-cell arithmetic.

\textbf{Bayes lower endpoint.}
Fix \(c\) and a maximal pattern in \(\mathcal S_c\). Atoms in the same cell \(\mathcal Q\) have identical revealed coordinates. Completing the square gives, for every action \(a\),
\[
\sum_{h\in\mathcal Q}q_h(\tau_h-a)^2
=\sum_{h\in\mathcal Q}q_h(\tau_h-m_{\mathcal Q})^2
+\left(\sum_{h\in\mathcal Q}q_h\right)(a-m_{\mathcal Q})^2.
\tag{6}
\]
Consequently the posterior mean \(m_{\mathcal Q}\) is the Bayes action and the minimized orbitwise risk is exactly \(B_c(\Pi_{b,s}^{\star})\). Orbit-constant weights and group closure give this same risk for every member of \(\mathcal S_c\).

For an arbitrary binary assignment, its revealed target-coordinate set is compatible and has a realizable maximal extension by \(\mathrm{lem{:}published\text{-}conflict\text{-}calibration}\). The extension contains all coordinates originally revealed. Conditional squared-error Bayes risk is conditional variance, and the law of total variance shows that conditioning on weakly more coordinates cannot increase its expectation. Thus the risk after the original assignment is at least the risk after its maximal extension, which is at least
\[
\min_{0\leq c\leq b}B_c(\Pi_{b,s}^{\star}).
\tag{7}
\]
Therefore every assignment law and every measurable estimator has Bayes risk under \(\Pi_{b,s}^{\star}\) at least (7). Worst-case risk over \(\Theta_{b,s}\) dominates Bayes risk under a prior supported on \(\Theta_{b,s}\), and hence
\[
\min_{0\leq c\leq b}B_c(\Pi_{b,s}^{\star})
\leq R_2(G_{b,s}).
\tag{8}
\]
Combining (2) and (8) proves the computable sandwich. If the finite calculation verifies \(\min_cB_c(\Pi_{b,s}^{\star})=V_{b,s}\), weak duality is tight at both ends: the purified procedure is minimax among all measurable procedures and the finite prior is least favorable for that verified instance. Without this endpoint equality, the argument establishes only the stated lower and upper certificates. \(\square\)
\end{proof}
\par\smallskip\noindent \textit{Justification.} This theorem gives an exact invariant-linear value, a realizable unrestricted upper action, and an auditable finite Bayes lower action on equal block cliques. \textit{Gap.} Published work does not contain the three-mode block-clique SDP or its finite posterior-cell lower certificate; the proof does not establish equality of the unrestricted lower and upper endpoints for every \(b,s\). \textit{Consumer.} CauseIS can benchmark against a computable certified interval and against an exact unrestricted value only when the endpoint check closes.

\begin{proposition}[prop:complete-clique-calibration]\label{prop:complete-clique-calibration}
For \(b=1\) and every \(s\ge2\), the SDP has value \(V_{1,s}=1\). The all-control design with estimator \(-s^{-1}\sum_{i\in V_n}u_{i,0}\) attains risk one. The uniform four-atom prior on layer-constant vectors \(u_{i,0}=x_0\), \(u_{i,1}=x_1\), \((x_0,x_1)\in\{-1,1\}^2\), is atomwise feasible and has posterior risk one both for the all-control orbit and for every singleton-treated orbit.
\end{proposition}
\begin{proof}
\textbf{Proof.}
Set \(b=1\), so \(n=s\). We first prove an SDP lower bound without assuming rank-one conditional moments. The constant-mode constraint implies
\[
\lambda_0+\lambda_1\geq\operatorname{tr}K_0(W).
\tag{1}
\]
Specializing the formula for \(K_0(W)\) to \(b=1\) gives
\[
\begin{aligned}
\operatorname{tr}K_0(W)
&=2w^0_{00}-2w^0_{01}+w^0_{11}
  +2w^1_{00}+2sw^1_{02}+s^2w^1_{22}\\
&=w^0_{00}+(1,-1,0)W_0(1,-1,0)^{\mathsf T}
  +w^1_{00}+(1,0,s)W_1(1,0,s)^{\mathsf T}\\
&\geq w^0_{00}+w^1_{00}=1.
\end{aligned}
\tag{2}
\]
The inequality follows from \(W_0,W_1\succeq0\), and the last equality is the SDP normalization. Thus every feasible point has objective at least one.

For the reverse inequality, take \(p_1=1\) and \(D_1=-1/s\); the unused value of \(A_1\) is arbitrary. The scaled constant-mode coefficient vector is \((0,-1)\), and there are no between-block or treated within-block modes. Hence \(K_0=\operatorname{diag}(0,1)\), so \((\lambda_0,\lambda_1)=(0,1)\) is feasible. Therefore
\[
V_{1,s}=1.
\tag{3}
\]

The corresponding all-control estimator is \(-s^{-1}\sum_i u_{i,0}\). Its error is
\[
-\frac1s\sum_i u_{i,0}
-\frac1s\sum_i(u_{i,1}-u_{i,0})
=-\frac1s\sum_i u_{i,1}.
\tag{4}
\]
By Cauchy--Schwarz and the treated-layer restriction,
\[
\left(\frac1s\sum_i u_{i,1}\right)^2
\leq\frac1s\sum_i u_{i,1}^2\leq1.
\tag{5}
\]
Equality holds when the treated layer is constant with value one, so the worst-case risk is exactly one.

For the lower certificate, let \(x_0,x_1\) be independent Rademacher variables and set \(u_{i,0}=x_0\) and \(u_{i,1}=x_1\) for every \(i\). Each of the four atoms satisfies both normalized layer constraints with equality. Under the all-control orbit the observation reveals \(x_0\), so
\[
\operatorname{Var}(x_1-x_0\mid x_0)=\operatorname{Var}(x_1)=1.
\tag{6}
\]
Under every singleton-treated orbit the observation reveals \(x_1\), so
\[
\operatorname{Var}(x_1-x_0\mid x_1)=\operatorname{Var}(x_0)=1.
\tag{7}
\]
The posterior mean is the squared-error Bayes action, proving posterior risk one in both orbit types. Any other binary assignment on the single clique reveals neither target sign and hence has posterior variance two; assigning the same non-target outcomes to all four atoms prevents those outcomes from conveying additional information. Thus every design--estimator pair has Bayes risk at least one. Together with (3)--(5), this also verifies \(R_2(G_{1,s})=V_{1,s}=1\). \(\square\)
\end{proof}
\par\smallskip\noindent \textit{Justification.} The complete clique gives a positive-risk exact example checking the SDP normalization, realizable upper coefficients, and a four-atom all-estimator Bayes lower certificate. \textit{Gap.} Published order bounds do not provide the exact finite-instance equality \(R_2(G_{1,s})=V_{1,s}=1\). \textit{Consumer.} The identity is a nondegenerate unit test for implementations of the block-clique SDP and posterior-risk calculations.

\begin{openendedquestion}[oeq:general-graph-finite-saddle-certificate]\label{oeq:general-graph-finite-saddle-certificate}
For an arbitrary DTE conflict graph, can the orbit-aware saddle column-generation handle produce in randomized polynomial time a finite prior lower certificate and a realizable design-estimator pair whose certified risks differ by at most a universal constant factor, with both supports polynomial in \(n\)?
\end{openendedquestion}
\textbf{Established partial result.}
For any nonempty ordered pair \(A,B\subseteq V_n\) satisfying
\[
B\subseteq\bigcap_{i\in A}\widetilde N_i(G),
\]
the proposition \(\mathrm{prop{:}balanced\text{-}biclique\text{-}four\text{-}atom\text{-}certificate}\) constructs a uniform four-atom prior in the unchanged class \(\mathcal M_2(G)\).  Its target layers are
\[
Y_i(\mathbf 0)=\sqrt{\frac{n}{|A|}}\,\xi_0\mathbf 1\{i\in A\},
\qquad
Y_i(\mathbf e_i)=\sqrt{\frac{n}{|B|}}\,\xi_1\mathbf 1\{i\in B\},
\]
where \(\xi_0\) and \(\xi_1\) are independent Rademacher signs.  Each atom satisfies arbitrary neighborhood interference, and the two separate layer moments equal
\[
\frac1n\sum_iY_i(\mathbf0)^2=1,
\qquad
\frac1n\sum_iY_i(\mathbf e_i)^2=1.
\]
Its direct treatment effect is
\[
\tau_{\mathrm{DTE}}^G
=
\sqrt{\frac{|B|}{n}}\,\xi_1
-
\sqrt{\frac{|A|}{n}}\,\xi_0.
\]
If an assignment reveals \(\xi_0\) through an all-control exposure at a vertex of \(A\), the containment condition forces every vertex of \(B\) to be assigned control, so no isolated-treatment exposure in \(B\) can reveal \(\xi_1\).  Conversely, revealing \(\xi_1\) requires a treated vertex in \(B\).  Hence no binary assignment reveals both signs.  Conditional on every realized assignment and observed outcome vector, at least one sign remains independent and posterior-unrevealed.  Completion of the square therefore gives, for every measurable estimator,
\[
\mathbb E\!\left[
\left\{\widehat\tau-\tau_{\mathrm{DTE}}^G\right\}^2
\,middle|\,z,Y(z)
\right]
\geq
\frac{\min\{|A|,|B|\}}{n}.
\]
Integrating over an arbitrary assignment law, using that worst-case risk dominates Bayes risk, and maximizing over feasible pairs proves the directly checkable lower certificate
\[
R_2(G)\geq\frac{\beta_{\mathrm{DTE}}(G)}{n}.
\]
The certificate has four atoms once a feasible pair is supplied.  Every clique \(Q\) supplies the pair \(A=B=Q\), so
\[
R_2(G)\geq\frac{\beta_{\mathrm{DTE}}(G)}n
\geq\frac{\omega(G)}n.
\]
For an equal block-clique graph, \(\beta_{\mathrm{DTE}}(G_{b,s})=s\), and the established computable sandwich strengthens to
\[
\max\left\{
\frac1b,
\min_{0\leq c\leq b}B_c(\Pi_{b,s}^{\star})
\right\}
\leq R_2(G_{b,s})\leq V_{b,s}.
\]
For \(G_{n,m}=K_{m,m}\sqcup(n-2m)K_1\), the two shores form a feasible ordered pair when \(m\geq1\), giving \(R_2(G_{n,m})\geq m/n\).  The established spectral upper functional gives \(R_2(G_{n,m})\lesssim(m+1)/n\); consequently \(R_2(G_{n,m})=\Theta(m/n)\) uniformly for \(m\geq1\).  Thus the lower-certificate side is constant-support and exactly auditable for a supplied feasible pair, but no randomized polynomial-time general-graph upper action or universal-factor two-sided certificate has been proved.
\par\smallskip\noindent \textit{Justification.} The balanced-biclique proposition establishes the constant-support lower certificate \(R_2(G)\geq\beta_{\mathrm{DTE}}(G)/n\), including the rescaled-clique corollary and complete-bipartite calibration. It does not produce the general upper action, so the two-sided polynomial-time saddle question remains open. \textit{Gap.} The missing step is a polynomial-time general-graph separation and realizability argument that outputs polynomial support on both sides and certifies a universal factor against the unrestricted all-estimator value. \textit{Consumer.} The proved lower certificate already refines CauseIS benchmarking on networks with large cross-layer supports; resolving the OEQ would additionally certify approximation ratios for CauseIS and Conflict Graph Design on arbitrary known networks.

\begin{lemma}[lem:published-network-risk-monotonicity]\label{lem:published-network-risk-monotonicity}
For the published \(q=2\) direct-treatment-effect problem, spanning-subgraph inclusion is risk monotone: if \(H_{\mathcal P}\subseteq G\) on the same vertex set, then
\[
R_2(H_{\mathcal P})\leq R_2(G).
\]
\end{lemma}
\par\smallskip\noindent \textit{Justification.} This cited leaf supplies the exact direction needed to transfer a lower bound from a spanning clique subgraph to the original graph. \textit{Gap.} The result is published substrate, located at Section 2.2, Proposition 2.1 and proved in Appendix A.2; it is not claimed as a contribution of this note. \textit{Consumer.} The clique-partition converse uses it after computing the all-procedure Bayes risk on \(H_{\mathcal P}\).

\begin{lemma}[lem:published-conflict-graph-characterization]\label{lem:published-conflict-graph-characterization}
For the published \(q=2\) direct-treatment-effect problem on a known graph \(G\), the unrestricted network minimax risk equals the minimax risk of the maximal-independent-set game on the labelled conflict graph \(H_{\mathrm{DTE}}(G)\); in particular, \(R_2(G)\) is a function only of that labelled conflict graph.
\end{lemma}

\begin{lemma}[lem:published-dte-graph-functional-bounds]\label{lem:published-dte-graph-functional-bounds}
For every known interference graph \(G\), the published \(q=2\) DTE risk satisfies
\[
\max\left\{\frac1{\alpha(G)},\frac{\sqrt{d_{\mathrm{avg}}(G)}}n\right\}
\lesssim R_2(G)
\lesssim\frac{\lambda(G)+1}{n},
\]
where \(\alpha(G)\) is the independence number, \(d_{\mathrm{avg}}(G)\) is the average degree, and \(\lambda(G)\) is the adjacency spectral radius.
\end{lemma}
\par\smallskip\noindent \textit{Justification.} This cited leaf provides the published graph-summary baseline used in the dense-family rate comparison. \textit{Gap.} The formal \(q=2\) DTE result is Appendix D.1, Corollary D.1; the new clique-partition functional can be strictly stronger in rate on the displayed family. \textit{Consumer.} It calibrates the improvement and the available spectral upper order without asserting a new general upper certificate.

\begin{proposition}[prop:block-product-rademacher-certificate]\label{prop:block-product-rademacher-certificate}
For every \(b\geq 1\) and \(s\geq 2\), consider the uniform \(4^b\)-atom prior obtained by taking all \(\xi_{r,k}\), \(1\leq r\leq b\), \(k\in\{0,1\}\), to be independent Rademacher signs and setting \(u_{i,k}=\xi_{r,k}\) whenever \(i\in B_r\). Every atom belongs to \(\Theta_{b,s}\). Its posterior risk is exactly \(1/b\) after every maximal observation pattern in every orbit \(\mathcal S_c\), and its posterior risk is at least \(1/b\) after every binary assignment. Consequently,
\[
\max\left\{\frac1b,\min_{0\leq c\leq b}B_c(\Pi_{b,s}^{\star})\right\}
\leq R_2(G_{b,s})\leq V_{b,s}.
\]
\end{proposition}
\begin{proof}
\textbf{Proof.}
Apply \(\mathrm{prop{:}clique\text{-}partition\text{-}rademacher\text{-}converse}\) to \(G_{b,s}\) and its defining partition \(\mathcal B_{b,s}\). Since all \(b\) blocks have size \(s\) and \(n=bs\), its general lower endpoint specializes to
\[
\sum_{r=1}^b\left(\frac{|B_r|}{n}\right)^2
=b\left(\frac{s}{bs}\right)^2=\frac1b.
\tag{1}
\]
This already gives the all-design, all-measurable-estimator lower bound. We additionally verify the proposition's finite posterior claim exactly.

Let the \(2b\) signs \((\xi_{r,k})\) be mutually independent Rademacher variables and put \(u_{i,k}=\xi_{r,k}\) for \(i\in B_r\). There are \(4^b\) equally weighted atoms, and for either layer
\[
\frac1n\sum_{i\in V_n}u_{i,k}^2
=\frac1{bs}\sum_{r=1}^b\sum_{i\in B_r}\xi_{r,k}^2=1.
\tag{2}
\]
Thus every atom belongs to \(\Theta_{b,s}\). Extend it to a schedule by using the target signs on all-control and isolated-treatment local patterns and zero on every other local pattern. This construction satisfies arbitrary neighborhood interference and the two separate second-moment restrictions.

The target is
\[
\tau_{\mathrm{DTE}}(Y)=\frac1b\sum_{r=1}^b(\xi_{r,1}-\xi_{r,0}).
\tag{3}
\]
Fix a maximal pattern in \(\mathcal S_c\), with all-control block set \(C\). If \(r\in C\), the data reveal \(\xi_{r,0}\) and not \(\xi_{r,1}\); if \(r\notin C\), they reveal \(\xi_{r,1}\) and not \(\xi_{r,0}\). Exactly one independent Rademacher sign remains unobserved in each block. Hence
\[
\operatorname{Var}\!\left(\tau_{\mathrm{DTE}}(Y)\mid S_c,u_{S_c}\right)
=\frac1{b^2}\sum_{r=1}^b1=\frac1b.
\tag{4}
\]
Conditional completion of the square shows that the posterior mean is the squared-loss Bayes action, so (4) is exactly the posterior risk after every member of every orbit \(\mathcal S_c\).

For an arbitrary assignment, each block is all control, has exactly one treated unit, or has at least two treated units. It reveals respectively the control sign, the treated sign, or neither sign, and never both. The conditional target variance is therefore
\[
\frac{\#\{\text{unrevealed signs}\}}{b^2}\geq\frac b{b^2}=\frac1b.
\tag{5}
\]
The fixed non-target outcomes carry no further information. Equations (4)--(5) reproduce the exact finite posterior calculation, while (1) records its specialization from the arbitrary-graph clique-partition certificate.

Finally, \(\mathrm{thm{:}block\text{-}clique\text{-}exact\text{-}saddle}\) supplies
\[
\min_{0\leq c\leq b}B_c(\Pi_{b,s}^{\star})
\leq R_2(G_{b,s})\leq V_{b,s}.
\tag{6}
\]
Taking the maximum of the two independently valid lower endpoints in (1) and (6) proves the displayed sandwich. No equality of its endpoints is used. \(\square\)
\end{proof}
\par\smallskip\noindent \textit{Justification.} The equal-block result is the explicit specialization of the arbitrary clique-partition converse and retains its exact orbitwise posterior calculation. \textit{Gap.} It supplies the closed-form \(1/b\) unrestricted lower endpoint but does not identify unrestricted risk unless a separate upper endpoint matches it. \textit{Consumer.} CauseIS obtains a directly interpretable lower benchmark alongside the block-clique SDP upper certificate.

\begin{proposition}[prop:clique-partition-rademacher-converse]\label{prop:clique-partition-rademacher-converse}
Let \(G\) be any simple graph on \(V_n\), and let \(\mathcal P\) be any clique partition of \(G\). Put mutually independent Rademacher signs \(\xi_{B,0},\xi_{B,1}\) on every \(B\in\mathcal P\), and on the spanning clique subgraph \(H_{\mathcal P}\) set the two target layers equal to \(\xi_{B,k}\) for every \(i\in B\) and \(k\in\{0,1\}\). Every one of the \(4^{|\mathcal P|}\) atoms belongs to the two separate normalized Euclidean balls. No binary assignment on \(H_{\mathcal P}\) reveals both signs from any block. After every maximal target-exposure observation its conditional DTE variance is exactly
\[
\sum_{B\in\mathcal P}\left(\frac{|B|}{n}\right)^2,
\]
and after every binary assignment it is at least this value. Consequently,
\[
R_2(H_{\mathcal P})\geq
\sum_{B\in\mathcal P}\left(\frac{|B|}{n}\right)^2,
\qquad
R_2(G)\geq R_2(H_{\mathcal P}),
\]
and hence the finite, directly enumerable clique-partition certificate is
\[
R_2(G)\geq
\max_{\substack{\mathcal P\text{ a partition of }V_n\\
G[B]\text{ complete for every }B\in\mathcal P}}
\sum_{B\in\mathcal P}\left(\frac{|B|}{n}\right)^2.
\]
No polynomial-time claim for maximizing this certificate is made.
\end{proposition}
\begin{proof}
\textbf{Proof.}
Fix \(G\) and a clique partition \(\mathcal P\), and write
\[
a_B=\frac{|B|}{n},\qquad B\in\mathcal P.
\tag{1}
\]
For every sign array \(\xi=(\xi_{B,k}:B\in\mathcal P,k\in\{0,1\})\), define a complete schedule on \(H_{\mathcal P}\) as follows.  If \(i\in B\), put
\[
Y_i^{\xi}(z)=
\begin{cases}
\xi_{B,0},&z|_B\equiv0,\\
\xi_{B,1},&z_i=1\text{ and }z_j=0\text{ for all }j\in B\setminus\{i\},\\
0,&\text{otherwise}.
\end{cases}
\tag{2}
\]
By the definition of \(H_{\mathcal P}\), the closed neighborhood of \(i\) is exactly its block \(B\).  Hence (2) is a complete arbitrary-neighborhood-interference schedule.  It satisfies
\[
Y_i^{\xi}(\mathbf0)=\xi_{B,0},
\qquad
Y_i^{\xi}(\mathbf e_i)=\xi_{B,1}.
\tag{3}
\]
Because the blocks partition \(V_n\), both separate target-layer restrictions hold with equality:
\[
\frac1n\sum_{i\in V_n}Y_i^{\xi}(\mathbf0)^2
=\frac1n\sum_{B\in\mathcal P}|B|=1,
\qquad
\frac1n\sum_{i\in V_n}Y_i^{\xi}(\mathbf e_i)^2=1.
\tag{4}
\]
Thus all \(4^{|\mathcal P|}\) atoms belong to \(\mathcal M_2(H_{\mathcal P})\).  Give them the uniform prior, so the signs are mutually independent.  The DTE is
\[
\tau_{\mathrm{DTE}}^{H_{\mathcal P}}(Y^{\xi})
=\sum_{B\in\mathcal P}a_B(\xi_{B,1}-\xi_{B,0}).
\tag{5}
\]

Fix a binary assignment \(z\).  Within a block \(B\), if \(z|_B\equiv0\), the observations reveal \(\xi_{B,0}\) but not \(\xi_{B,1}\).  If exactly one unit is treated, its outcome reveals \(\xi_{B,1}\), whereas every other observed outcome in the block is the fixed value zero, so \(\xi_{B,0}\) is not revealed.  If at least two units are treated, all observed outcomes in the block are zero and neither sign is revealed.  This exhaustive split also covers singleton blocks.  Thus no assignment reveals both signs from any block.

All unrevealed signs remain conditionally independent Rademacher variables, and therefore
\[
\operatorname{Var}_{\xi}\!\left(
\tau_{\mathrm{DTE}}^{H_{\mathcal P}}(Y^{\xi})
\mid z,Y^{\xi}(z)
\right)
=\sum_{B\in\mathcal P}a_B^2N_B(z),
\tag{6}
\]
where \(N_B(z)\in\{1,2\}\) is the number of unrevealed signs in block \(B\).  Hence (6) is at least \(\sum_Ba_B^2\).  By the cited conflict-graph characterization, a maximal target-exposure observation is a maximal compatible exposure pattern.  Within each disjoint clique block it contains either the all-control exposure or one isolated-treatment exposure; otherwise another compatible target exposure could be added.  It consequently reveals exactly one sign per block, so \(N_B(z)=1\) for every \(B\) and (6) equals \(\sum_Ba_B^2\).

For squared loss and every measurable estimator value \(t\), conditional completion of the square gives
\[
\mathbb E_{\xi}[(\tau-t)^2\mid z,Y^{\xi}(z)]
\geq\operatorname{Var}_{\xi}(\tau\mid z,Y^{\xi}(z))
\geq\sum_{B\in\mathcal P}a_B^2.
\tag{7}
\]
Integrating over an arbitrary assignment law, then using that worst-case risk dominates finite-prior Bayes risk, proves
\[
R_2(H_{\mathcal P})\geq
\sum_{B\in\mathcal P}\left(\frac{|B|}{n}\right)^2.
\tag{8}
\]
The subgraph \(H_{\mathcal P}\subseteq G\) is spanning, so the cited network-risk monotonicity lemma gives
\[
R_2(G)\geq R_2(H_{\mathcal P}).
\tag{9}
\]
There are finitely many partitions of the finite set \(V_n\).  Maximizing (8)--(9) over exactly those partitions whose blocks induce cliques gives the displayed finite enumeration.  No step supplies a polynomial-time maximization algorithm. \(\square\)
\end{proof}
\par\smallskip\noindent \textit{Justification.} The proposition remains a finite blockwise Bayes certificate and now serves as a dominated but directly enumerable comparison point. \textit{Gap.} The published lower functional does not state this weighted partition posterior variance, although the stronger balanced-biclique certificate dominates it through the clique corollary. \textit{Consumer.} CauseIS can retain the partition value as an easily audited lower benchmark on networks with an available clique partition.

\begin{proposition}[prop:rescaled-clique-four-atom-certificate]\label{prop:rescaled-clique-four-atom-certificate}
Let \(G\) be any simple graph on the nonempty finite population \(V_n\).  For every supplied nonempty clique \(Q\subseteq V_n\), writing \(q=|Q|\) and \(a=\sqrt{n/q}\), the uniform four-atom prior with independent Rademacher signs \(\xi_0,\xi_1\), target outcomes \(Y_i(\mathbf 0)=a\xi_0\) and \(Y_i(\mathbf e_i)=a\xi_1\) for \(i\in Q\), and zero target outcomes off \(Q\), extends to complete schedules in \(\mathcal M_2(G)\).  Every binary assignment reveals at most one of \(\xi_0,\xi_1\), and consequently every assignment law and measurable estimator satisfies
\[
R_2(G)\geq \frac{q}{n}.
\]
Maximizing over cliques gives
\[
R_2(G)\geq\frac{\omega(G)}{n}.
\]
Moreover, for every clique partition \(\mathcal P\) of \(G\),
\[
\sum_{B\in\mathcal P}\left(\frac{|B|}{n}\right)^2
\leq \frac{\omega(G)}{n},
\]
so the four-atom maximum-clique certificate dominates the clique-partition certificate.  A supplied clique makes \(q/n\) directly computable, but no polynomial-time maximum-clique algorithm is asserted.  On \(G_{b,s}\), \(\omega(G_{b,s})=s\) and the bound is \(1/b\).
\end{proposition}
\begin{proof}
\textbf{Proof.}
Fix a nonempty clique \(Q\) and put \(q=|Q|\).  For every \(i\in Q\), clique completeness gives \(Q\subseteq\widetilde N_i(G)\).  Hence the ordered pair \((A,B)=(Q,Q)\) is feasible in the balanced-biclique certificate.  Its two amplitudes both equal \(a=\sqrt{n/q}\), and the complete schedules constructed there have exactly the target outcomes displayed in this proposition and zero target outcomes off \(Q\).  The atomwise support calculation, the observation incompatibility, and the posterior-variance calculation in that certificate therefore give
\[
R_2(G)\geq\frac{q}{n}.
\tag{1}
\]
Because the finite nonempty graph has a maximum clique, maximizing (1) gives \(R_2(G)\geq\omega(G)/n\).

For any clique partition \(\mathcal P\), every block has size at most \(\omega(G)\), and the blocks have total size \(n\).  Thus
\[
\sum_{B\in\mathcal P}\left(\frac{|B|}{n}\right)^2
\leq\frac{\omega(G)}{n^2}\sum_{B\in\mathcal P}|B|
=\frac{\omega(G)}n.
\tag{2}
\]
The clique-partition proposition proves that the left-hand side of (2) is itself a valid all-procedure lower certificate, so (2) proves the asserted domination.

Finally, every clique of \(G_{b,s}\) lies within one component block, and each block is a clique of size \(s\).  Therefore \(\omega(G_{b,s})=s\).  Since \(n=bs\), (1) gives \(s/n=1/b\). \(\square\)
\end{proof}
\par\smallskip\noindent \textit{Justification.} The useful supplied-clique witness is retained as the symmetric special case \(A=B=Q\) of the stronger cross-layer certificate. \textit{Gap.} Clique support alone misses complete-bipartite cross-layer collisions, but remains a simple directly verifiable corollary. \textit{Consumer.} CauseIS can compare its risk certificate with \(|Q|/n\) whenever a dense clique witness is already available.

\begin{proposition}[prop:polynomial-clique-isolates-four-atom-calibration]\label{prop:polynomial-clique-isolates-four-atom-calibration}
Fix \(\rho\in(0,1)\) and \(\gamma\in(0,1]\).  For every \(n\geq1\), let \(m_n=\lfloor\rho n^{\gamma}\rfloor\), and let \(G_n\) be the disjoint union of \(K_{m_n}\) and \(n-m_n\) isolated vertices, interpreting \(K_0\) as the empty graph.  Then \(0\leq m_n\leq n-1\),
\[
\omega(G_n)=\max\{1,m_n\},
\qquad
R_2(G_n)\geq\frac{\max\{1,m_n\}}n,
\]
and the exact graph quantities are
\[
\alpha(G_n)=
\begin{cases}
n,&m_n=0,\\
n-m_n+1,&m_n\geq1,
\end{cases}
\qquad
d_{\mathrm{avg}}(G_n)=\frac{m_n(m_n-1)}n,
\qquad
\lambda(G_n)=
\begin{cases}
0,&m_n=0,\\
m_n-1,&m_n\geq1.
\end{cases}
\]
The four-atom clique lower and the cited spectral upper functional have the same order for every \(\gamma\in(0,1]\):
\[
R_2(G_n)=\Theta(n^{\gamma-1}),
\qquad
\frac{\lambda(G_n)+1}{n}
=\frac{\max\{1,m_n\}}n
=\Theta(n^{\gamma-1}).
\]
The displayed published lower functional has order
\[
\max\left\{\frac1{\alpha(G_n)},
\frac{\sqrt{d_{\mathrm{avg}}(G_n)}}n\right\}
=\Theta\!\left(n^{-1}\vee n^{\gamma-3/2}\right),
\]
so the new lower improves that displayed functional by
\[
\Theta\!\left(n^{\min\{\gamma,1/2\}}\right)
\]
for every \(\gamma\in(0,1]\).  The older clique-plus-singletons partition certificate \((m_n^2+n-m_n)/n^2\) remains valid but is never larger than \(\max\{1,m_n\}/n\).  No efficient maximum-clique algorithm or new arbitrary-graph upper design is asserted.
\end{proposition}
\begin{proof}
\textbf{Proof.}
Write \(m=m_n\).  Since \(0<\rho<1\), \(0<\gamma\leq1\), and \(n\geq1\),
\[
0\leq m=\lfloor\rho n^{\gamma}\rfloor<n,
\tag{1}
\]
so \(0\leq m\leq n-1\).  If \(m=0\), all vertices are isolated and the clique number is one because \(V_n\) is nonempty.  If \(m=1\), every component is again a singleton.  If \(m\geq2\), the \(K_m\) component is a clique of size \(m\), and no clique can meet two distinct components.  Therefore
\[
\omega(G_n)=\max\{1,m\}.
\tag{2}
\]
Applying \(\text{prop:rescaled-clique-four-atom-certificate}\) to a maximum clique gives the finite lower certificate
\[
R_2(G_n)\geq\frac{\max\{1,m\}}n.
\tag{3}
\]

The independence number is \(n\) when \(m=0\).  When \(m\geq1\), an independent set may contain all \(n-m\) isolated vertices and at most one vertex of \(K_m\), giving \(\alpha(G_n)=n-m+1\).  The degree-sum identity and the block-diagonal adjacency matrix give
\[
d_{\mathrm{avg}}(G_n)=\frac{2\binom m2}{n}=\frac{m(m-1)}n,
\tag{4}
\]
and \(\lambda(G_n)=0\) for \(m=0\), while \(\lambda(G_n)=m-1\) for \(m\geq1\).  In particular,
\[
\frac{\lambda(G_n)+1}{n}=\frac{\max\{1,m\}}n
\tag{5}
\]
at every finite \(n\).

Because \(\rho n^{\gamma}\to\infty\), for all sufficiently large \(n\),
\[
\frac{\rho}{2}n^{\gamma}\leq m\leq\rho n^{\gamma}
\quad\text{and}\quad m\geq2.
\tag{6}
\]
Equations (3), (5), and (6), together with the upper inequality in \(\text{lem:published-dte-graph-functional-bounds}\), give
\[
R_2(G_n)=\Theta(n^{\gamma-1}).
\tag{7}
\]
This uses the published spectral upper comparator; it does not construct a new arbitrary-graph design.

For the displayed published lower functional, the restriction \(\rho<1\) implies \(n-m=\Theta(n)\), including when \(\gamma=1\).  Hence
\[
\frac1{\alpha(G_n)}=\Theta(n^{-1}),
\qquad
\frac{\sqrt{d_{\mathrm{avg}}(G_n)}}n
=\frac{\sqrt{m(m-1)}}{n^{3/2}}
=\Theta(n^{\gamma-3/2}).
\tag{8}
\]
Thus that displayed lower functional has order \(n^{-1}\vee n^{\gamma-3/2}\).  Dividing the new order \(n^{\gamma-1}\) by this functional gives
\[
\frac{n^{\gamma-1}}{n^{-1}\vee n^{\gamma-3/2}}
=
\begin{cases}
\Theta(n^{\gamma}),&0<\gamma\leq1/2,\\
\Theta(n^{1/2}),&1/2\leq\gamma\leq1,
\end{cases}
=\Theta\!\left(n^{\min\{\gamma,1/2\}}\right).
\tag{9}
\]
The two cases agree at \(\gamma=1/2\), so there is no onset threshold for improvement: the factor diverges for every fixed \(\gamma>0\).

For completeness, the older partition value is
\[
L_n=\frac{m^2+n-m}{n^2}.
\tag{10}
\]
If \(m=0\), then \(L_n=1/n=\max\{1,m\}/n\).  If \(m\geq1\), direct subtraction gives
\[
\frac mn-L_n
=\frac{(m-1)(n-m)}{n^2}\geq0,
\tag{11}
\]
using (1).  This proves the asserted finite domination and completes all boundary cases. \(\square\)
\end{proof}
\par\smallskip\noindent \textit{Justification.} This calibration proves that the new lower and cited spectral upper have order \(n^{\gamma-1}\) for every positive clique-size exponent. \textit{Gap.} The displayed published lower functional is smaller by a diverging factor throughout the entire polynomial regime, not only beyond a \(\gamma=1/2\) onset. \textit{Consumer.} CauseIS benchmarking can use the matched minimax order across all polynomial clique sizes without reporting a spurious intermediate upper gap.

\begin{proposition}[prop:dense-clique-isolates-four-atom-calibration]\label{prop:dense-clique-isolates-four-atom-calibration}
Fix \(\rho\in(0,1)\), let \(m_n=\lfloor\rho n\rfloor\), and let \(G_n\) be the disjoint union of \(K_{m_n}\) and \(n-m_n\) isolated vertices.  The rescaled four-atom clique certificate and the cited spectral upper functional give
\[
R_2(G_n)\geq\frac{\max\{1,m_n\}}n\longrightarrow\rho,
\qquad
R_2(G_n)=\Theta(1).
\]
The older clique-partition certificate remains valid and converges to \(\rho^2\), but it is pointwise no larger than the four-atom certificate.  The two terms in the displayed published general DTE lower functional have orders \(n^{-1}\) and \(n^{-1/2}\), whereas the new lower and the cited spectral upper functional are both of order one.  This is a lower-certificate calibration for the published \(q=2\) DTE problem; it does not supply a new arbitrary-graph upper design.
\end{proposition}
\begin{proof}
\textbf{Proof.}
Apply \(\text{prop:polynomial-clique-isolates-four-atom-calibration}\) with \(\gamma=1\).  Since \(m_n/n\to\rho>0\), eventually \(m_n\geq1\), and its exact finite certificate becomes
\[
R_2(G_n)\geq\frac{m_n}{n}\longrightarrow\rho.
\tag{1}
\]
The same proposition and the cited spectral upper inequality give \(R_2(G_n)=\Theta(1)\).  The older partition certificate is
\[
\frac{m_n^2+n-m_n}{n^2}
=\left(\frac{m_n}{n}\right)^2+\frac{n-m_n}{n^2}
\longrightarrow\rho^2,
\tag{2}
\]
and its pointwise domination follows from equation (11) in that proposition.  Finally, its exact graph formulas give
\[
\frac1{\alpha(G_n)}=\Theta(n^{-1}),
\qquad
\frac{\sqrt{d_{\mathrm{avg}}(G_n)}}n=\Theta(n^{-1/2}),
\qquad
\frac{\lambda(G_n)+1}{n}=\Theta(1).
\tag{3}
\]
These are respectively the two terms in the displayed published lower functional and its spectral upper functional.  Equations (1)--(3) prove every assertion. \(\square\)
\end{proof}
\par\smallskip\noindent \textit{Justification.} The fixed-density specialization exposes the stronger lower limit \(\rho\), while recording the older partition limit \(\rho^2\) only as a dominated certificate. \textit{Gap.} The anchor's displayed lower functional is only of order \(n^{-1/2}\) on this family, despite order-one lower and upper certificates. \textit{Consumer.} CauseIS experiments with one dense community and isolated units should benchmark against an order-one minimax obstruction.

\begin{proposition}[prop:balanced-biclique-four-atom-certificate]\label{prop:balanced-biclique-four-atom-certificate}
Let \(G\) be any simple graph on the nonempty finite population \(V_n\).  Suppose that \(A,B\subseteq V_n\) are nonempty and
\[
B\subseteq\bigcap_{i\in A}\widetilde N_i(G).
\]
Let \(a_0=\sqrt{n/|A|}\) and \(a_1=\sqrt{n/|B|}\).  There is a uniform four-atom prior on complete schedules in the unchanged published class \(\mathcal M_2(G)\) whose target layers are
\[
Y_i(\mathbf 0)=a_0\xi_0\mathbf 1\{i\in A\},
\qquad
Y_i(\mathbf e_i)=a_1\xi_1\mathbf 1\{i\in B\},
\]
where \(\xi_0,\xi_1\) are independent Rademacher signs.  Every atom saturates each separate layer moment constraint exactly, and no binary assignment reveals both signs.  Consequently every assignment law and measurable estimator has worst-case mean squared error at least
\[
\frac{\min\{|A|,|B|\}}{n},
\]
so
\[
R_2(G)\geq\frac{\beta_{\mathrm{DTE}}(G)}{n}.
\]
Moreover, \(1\leq\beta_{\mathrm{DTE}}(G)\leq n\) and \(\beta_{\mathrm{DTE}}(G)\geq\omega(G)\).
\end{proposition}
\begin{proof}
\textbf{Proof.}
Fix a feasible ordered pair \((A,B)\).  Nonemptiness gives \(1\leq |A|,|B|\leq n\), so
\[
a_0=\sqrt{\frac{n}{|A|}},
\qquad
a_1=\sqrt{\frac{n}{|B|}}
\tag{1}
\]
are finite and positive.  Let \(\xi_0,\xi_1\) be independent Rademacher signs.  For each of their four values define a complete schedule, for every \(i\in V_n\) and \(z\in\mathcal Z_n\), by
\[
Y_i^{\xi}(z)=
\begin{cases}
a_0\xi_0,
& i\in A\text{ and }z|_{\widetilde N_i(G)}\equiv0,\\
a_1\xi_1,
& i\in B,\ z_i=1,\text{ and }z_j=0
   \text{ for all }j\in\widetilde N_i(G)\setminus\{i\},\\
0,&\text{otherwise}.
\end{cases}
\tag{2}
\]
The first two cases are disjoint because they respectively require \(z_i=0\) and \(z_i=1\).  Formula (2) assigns an outcome for every unit and every binary assignment, including every non-target exposure, and it depends on \(z\) only through \(z|_{\widetilde N_i(G)}\).  Thus every atom satisfies arbitrary neighborhood interference in the definition of \(\mathcal M_2(G)\).  At the two target assignments,
\[
Y_i^{\xi}(\mathbf 0)=a_0\xi_0\mathbf 1\{i\in A\},
\qquad
Y_i^{\xi}(\mathbf e_i)=a_1\xi_1\mathbf 1\{i\in B\}.
\tag{3}
\]
Consequently each of the two separate moment restrictions is saturated:
\[
\frac1n\sum_{i\in V_n}Y_i^{\xi}(\mathbf0)^2
=\frac{|A|a_0^2}{n}=1,
\qquad
\frac1n\sum_{i\in V_n}Y_i^{\xi}(\mathbf e_i)^2
=\frac{|B|a_1^2}{n}=1.
\tag{4}
\]
Equations (2)--(4) prove atomwise membership in the exact published class, without any restriction on non-target outcomes beyond the required neighborhood dependence.

The DTE of an atom is
\[
\tau_{\mathrm{DTE}}^G(Y^{\xi})
=\frac{|B|a_1}{n}\xi_1-\frac{|A|a_0}{n}\xi_0
=\sqrt{\frac{|B|}{n}}\,\xi_1
-\sqrt{\frac{|A|}{n}}\,\xi_0.
\tag{5}
\]
We now use the cross-layer containment.  Fix any assignment \(z\).  If an observed outcome depends on \(\xi_0\), the first line of (2) holds for some \(i\in A\).  Every vertex in \(B\) then has assignment zero because \(B\subseteq\widetilde N_i(G)\).  If an observed outcome depends on \(\xi_1\), the second line holds for some \(j\in B\), and hence \(z_j=1\).  These two events cannot occur simultaneously.  All other observed outcomes equal the fixed value zero.  Therefore the observation \((z,Y^{\xi}(z))\) depends on at most one of the two independent signs.

At least one sign remains an independent Rademacher variable after conditioning on the observation.  From (5), its squared coefficient is either \(|A|/n\) or \(|B|/n\); hence
\[
\operatorname{Var}_{\xi}\!\left(
\tau_{\mathrm{DTE}}^G(Y^{\xi})
\mid z,Y^{\xi}(z)
\right)
\geq\frac{\min\{|A|,|B|\}}{n}.
\tag{6}
\]
For every measurable estimator value \(t=t(z,Y^{\xi}(z))\), conditional completion of the square gives
\[
\mathbb E_{\xi}\!\left[
\{t-\tau_{\mathrm{DTE}}^G(Y^{\xi})\}^2
\mid z,Y^{\xi}(z)
\right]
=\operatorname{Var}_{\xi}(\tau_{\mathrm{DTE}}^G(Y^{\xi})\mid z,Y^{\xi}(z))
+\{t-\mathbb E_{\xi}[\tau_{\mathrm{DTE}}^G(Y^{\xi})\mid z,Y^{\xi}(z)]\}^2
\geq\frac{\min\{|A|,|B|\}}{n}.
\tag{7}
\]
Integrating (7) over any assignment law, using that worst-case risk dominates the average risk under this four-atom prior, and then taking the joint infimum over all assignment laws and measurable estimators proves the pairwise lower bound.  The population is finite, and the feasible set is nonempty because \(A=B=\{i\}\) is feasible for every \(i\in V_n\).  Maximizing the pairwise bound therefore proves
\[
R_2(G)\geq\frac{\beta_{\mathrm{DTE}}(G)}{n},
\qquad
1\leq\beta_{\mathrm{DTE}}(G)\leq n.
\tag{8}
\]
Finally, if \(Q\) is a nonempty clique, then \(Q\subseteq\widetilde N_i(G)\) for every \(i\in Q\).  Thus \((A,B)=(Q,Q)\) is feasible and \(\beta_{\mathrm{DTE}}(G)\geq |Q|\).  Taking a maximum clique proves \(\beta_{\mathrm{DTE}}(G)\geq\omega(G)\). \(\square\)
\end{proof}
\par\smallskip\noindent \textit{Justification.} This proposition is the strongest arbitrary-graph lower certificate delivered by the round and is uniform over all designs and measurable estimators. \textit{Gap.} Published graph functionals and MIS proxies do not expose this exact four-atom posterior obstruction on unequal cross-layer supports. \textit{Consumer.} CauseIS can use a supplied feasible pair as a rigorous lower benchmark for its MIS-based DTE design.

\begin{proposition}[prop:rescaled-clique-isolates-rate-calibration]\label{prop:rescaled-clique-isolates-rate-calibration}
Fix \(\rho\in(0,1/2)\) and \(\gamma\in(0,1]\).  For \(n\geq1\), let \(m_n=\lfloor\rho n^{\gamma}\rfloor\), and let \(G_n\) be the disjoint union of \(K_{m_n,m_n}\) and \(n-2m_n\) isolated vertices, with \(K_{0,0}\) interpreted as empty.  Then \(0\leq2m_n\leq n\).  If \(m_n\geq1\), orient the two shores as the control support \(A\) and treated support \(B\).  This pair satisfies
\[
B\subseteq\bigcap_{i\in A}\widetilde N_i(G_n),
\]
and hence
\[
\beta_{\mathrm{DTE}}(G_n)\geq
\begin{cases}
1,&m_n=0,\\
\max\{2,m_n\},&m_n\geq1.
\end{cases}
\]
Moreover,
\[
\omega(G_n)=
\begin{cases}
1,&m_n=0,\\
2,&m_n\geq1,
\end{cases}
\qquad
\lambda(G_n)=m_n.
\]
The balanced-biclique lower certificate and the cited spectral upper functional therefore match uniformly up to universal constants:
\[
R_2(G_n)\asymp\frac{m_n+1}{n},
\]
and in particular \(R_2(G_n)=\Theta(n^{\gamma-1})\) as \(n\to\infty\).  At the finite boundaries \(m_n=0\) and \(m_n=1\), the certified lower orders are respectively \(1/n\) and \(2/n\).  When \(m_n\to\infty\), the clique certificate uses only \(\omega(G_n)=2\) and gives \(2/n\), whereas the cross-layer certificate gives at least \(m_n/n\).  No claim is made that maximizing \(\beta_{\mathrm{DTE}}(G)\) is polynomial time.
\end{proposition}
\begin{proof}
\textbf{Proof.}
Because \(0<\gamma\leq1\), one has \(n^{\gamma}\leq n\) for \(n\geq1\).  Since \(\rho<1/2\),
\[
0\leq m_n\leq\rho n^{\gamma}<\frac n2.
\tag{1}
\]
As \(m_n\) is an integer, (1) implies \(0\leq2m_n\leq n\), so the graph is well defined.

Suppose first that \(m_n\geq1\), and denote the two nonempty shores of \(K_{m_n,m_n}\) by \(L\) and \(R\).  For every \(i\in L\),
\[
\widetilde N_i(G_n)=\{i\}\cup R,
\tag{2}
\]
because there are no edges within a shore or to the isolated vertices.  In particular,
\[
R\subseteq\bigcap_{i\in L}\widetilde N_i(G_n).
\tag{3}
\]
Thus the ordered cross-layer choice \(A=L\), \(B=R\) is feasible and gives \(\beta_{\mathrm{DTE}}(G_n)\geq m_n\).  The orientation matters: observing the control sign on \(L\) forces all of \(R\) to control, whereas observing an isolated-treatment sign on \(R\) treats a vertex of \(R\).  The general four-atom certificate proves the resulting lower risk \(m_n/n\).

If \(m_n=0\), the graph consists of \(n\) isolated vertices, and the universal singleton feasible pair gives \(\beta_{\mathrm{DTE}}(G_n)\geq1\).  The clique number is one in this case.  If \(m_n\geq1\), the bipartite component contains an edge and has no triangle, while isolates create no larger clique, so \(\omega(G_n)=2\).  Since the balanced-biclique parameter dominates clique number, these facts combine with (3) to give
\[
\beta_{\mathrm{DTE}}(G_n)\geq
\ell_n:=
\begin{cases}
1,&m_n=0,\\
\max\{2,m_n\},&m_n\geq1.
\end{cases}
\tag{4}
\]
The lower certificate yields \(R_2(G_n)\geq\ell_n/n\).

For \(m_n\geq1\), the adjacency matrix on the bipartite component is
\[
\begin{pmatrix}0&J_{m_n}\\J_{m_n}&0\end{pmatrix}.
\tag{5}
\]
The all-ones matrix \(J_{m_n}\) has one singular value \(m_n\) and all remaining singular values zero.  Hence (5) has eigenvalues \(m_n,-m_n\), and zero; adding isolated vertices adds only zero eigenvalues.  Thus \(\lambda(G_n)=m_n\) for \(m_n\geq1\), and it is also true for \(m_n=0\).  The cited published spectral upper functional supplies a universal constant \(C_{\mathrm{pub}}<\infty\) such that
\[
R_2(G_n)\leq C_{\mathrm{pub}}\frac{m_n+1}{n}.
\tag{6}
\]
For every nonnegative integer \(m_n\), (4) gives \(\ell_n\geq(m_n+1)/2\).  Combining this fact with (4) and (6) proves
\[
\frac{m_n+1}{2n}\leq R_2(G_n)
\leq C_{\mathrm{pub}}\frac{m_n+1}{n}.
\tag{7}
\]
This is the claimed uniform two-sided rate and verifies the cases \(m_n=0\) and \(m_n=1\) directly.  Finally, \(m_n\sim\rho n^{\gamma}\), so (7) gives \(R_2(G_n)=\Theta(n^{\gamma-1})\).  When \(m_n\geq1\), the clique-only certificate is exactly \(\omega(G_n)/n=2/n\); (3) gives the stronger \(m_n/n\) term once \(m_n>2\). \(\square\)
\end{proof}
\par\smallskip\noindent \textit{Justification.} The complete-bipartite-plus-isolates family proves that the new lower certificate can be rate-sharp where clique number stays bounded. \textit{Gap.} A clique-only lower is \(2/n\), whereas the shore certificate and cited spectral upper are both of order \(m_n/n\). \textit{Consumer.} CauseIS obtains a triangle-free stress test whose all-procedure benchmark scales with the cross-layer shore size.

\section{Interpretation}

The causal target is the fixed finite-population contrast \(\tau_{\mathrm{DTE}}^G(Y)\). The minimax risk \(R_2(G)\), the four-atom Bayes risk, the mode-SDP value \(V_{b,s}\), and the published spectral upper functional are decision-theoretic identifying or certification functionals for that target; none is the causal effect itself. A supplied balanced-biclique pair makes the lower certificate directly evaluable, but optimizing \(\beta_{\mathrm{DTE}}(G)\) is not claimed to be polynomial time.

\section*{Technical internal limitation (diagnostic only)}

The balanced-biclique lower certificate has four atoms once a feasible pair is supplied, but the paper gives no polynomial-time algorithm for maximizing \(\beta_{\mathrm{DTE}}(G)\). The equal block-clique posterior dictionary may be exponential in the symmetry-group size, and its lower endpoint need not equal the invariant-linear SDP endpoint without an explicit posterior-cell verification. No general-network realizable upper action, polynomial-support saddle, Unique-Games hardness reduction, or universal constant-factor approximation theorem is established.

\section{Honest scope}

Established results are the general supplied-pair lower bound \(R_2(G)\geq\min\{|A|,|B|\}/n\), its optimized certificate \(R_2(G)\geq\beta_{\mathrm{DTE}}(G)/n\), the clique and partition corollaries, the complete-bipartite rate calibration, and the equal block-clique invariant-linear SDP with a computable unrestricted sandwich. Unrestricted exact equality on the block-clique family is conditional on endpoint verification, except for the complete-clique case \(b=1\). The requested randomized polynomial-time two-sided constant-factor saddle on arbitrary DTE conflict graphs remains an open question; this paper neither proves the proposed \(d^{1/2+o(1)}\) frontier nor supplies its missing general upper certificate.

\begin{thebibliography}{99}

  \bibitem{KandirosHarshawSavje2026} Vardis Kandiros, Christopher Harshaw, and Fredrik Sävje. A Design-Based Minimax Theory for Network Experiments. arXiv:2608.04909v1, 2026. Section 2.2, Proposition 2.1 and Appendix A.2; Section 2.3, Proposition 2.2 and Appendix A.5; Appendix D.1, Corollary D.1.

  \bibitem{Gabler1990} Siegfried Gabler. Minimax Solutions in Sampling from Finite Populations. Lecture Notes in Statistics 64, Springer, 1990; electronic edition, 2012. Chapter 2, pages 20--24, and the quadratic parameter-space chapters. DOI: 10.1007/978-1-4612-3442-5.

  \bibitem{Rinott2009} Yosef Rinott. Some Decision-Theoretic Aspects of Finite Population Sampling. Handbook of Statistics 29, pages 523--558, 2009. DOI: 10.1016/S0169-7161(09)00241-7.

  \bibitem{KandirosPipisDaskalakisHarshaw2024} Vardis Kandiros, Charilaos Pipis, Constantinos Daskalakis, and Christopher Harshaw. The Conflict Graph Design: Estimating Causal Effects under Arbitrary Neighborhood Interference. arXiv:2411.10908v3, 2024.

  \bibitem{FatemiZheleva2023} Zahra Fatemi and Elena Zheleva. Minimizing Interference and Selection Bias in Network Experiment Design. Frontiers in Big Data 6:1128649, 2023.

\end{thebibliography}

\end{document}